\documentclass[a4paper, 12pt]{extarticle}


%%%%%%%%%%%%%%%%%%%%%%%%%%%%%%
% PACKAGES
%%%%%%%%%%%%%%%%%%%%%%%%%%%%%%

%!TEX encoding = UTF8

%%%%%%%%%%%%%%%%%%%%%%%%%%%%%%%%%
% PACKAGE IMPORTS
%%%%%%%%%%%%%%%%%%%%%%%%%%%%%%%%%


\usepackage[french]{babel}
\usepackage[T1]{fontenc}
\usepackage[utf8]{inputenc}

\usepackage[
	margin=2cm,
	right=2.2cm,
	footskip=1.5cm,
]{geometry}

%ams
\usepackage{amsmath,amsfonts,amsthm,amssymb,mathtools}

% for fancy headers
\usepackage{fancyhdr}

% for table of content stuff
\usepackage{titlesec}

\usepackage{bookmark}
\usepackage{enumitem}
\usepackage[most,many,breakable]{tcolorbox}
\usepackage{varwidth,etoolbox}
\usepackage[makeroom]{cancel}
\usepackage{xcolor}
\usepackage{multicol,array}
% algorithms
\usepackage[ruled,vlined,linesnumbered]{algorithm2e}

% delays answers
% answers are from last exercise if no label is provided
% delays exercices because of exercise bank (exercices.tex file in course folder)
%\usepackage[answerdelayed, exercisedelayed, lastexercise]{exercise}
\usepackage[answerdelayed, lastexercise]{exercise}

% for slanted sfrac
\usepackage{xfrac}
% for \uline which allows underline with line breaks
\usepackage[normalem]{ulem}

\usepackage{caption, subcaption}

% margins for points (and old (p. %page%) ).
\usepackage{marginnote}

\usepackage{makecell} % pour \thead 

%index
\usepackage[texindy]{imakeidx}
\makeindex[
	columns=2, 
	title=Index, 
	options= -L french,
	%intoc,
]
\indexsetup{level=\chapter}

%virgules
\usepackage{icomma}

% for striked out implies sign (\centernot\implies)
\usepackage{centernot}

% roman numerals for \section
\renewcommand{\thesection}{\Roman{section}} 

% tikz
\usepackage{tikz}

% package systeme pour les systèmes d'équations bien alignés
\usepackage{systeme}
% add \alpha \beta \gamma as variables
% put \\ as separator so commas don't return to line
\setsysteme{eq sep = \\, align={r,r}}
% example : 
% \systeme[sort={\alpha,\beta, \gamma}]{ \alpha + 2\beta = 4, \\ \beta = 8.} 

% for polynomial long division
\usepackage{polynom}
% example : 
% \polylongdiv[style=D]{x^3 - 1 + 2x - 2x^2}{x-1}


% date
\usepackage{advdate}


% tikz
\usepackage{tikz, pgfplots}
\usepackage{prerex}
\pgfplotsset{
	/pgf/number format/.cd,
	use comma,
	1000 sep={}
}

% for function graph
\usetikzlibrary{positioning}
\usetikzlibrary{shapes.geometric}
\usetikzlibrary{positioning}
\usetikzlibrary{shapes.misc}
\tikzset{
dot/.style = {circle, fill=#1, minimum size=5pt,
              inner sep=0pt, outer sep=0pt},
dot/.default = black % size of the circle diameter
}
\tikzset{cross/.style={cross out, draw=black, minimum size=2*(#1-\pgflinewidth), inner sep=0pt, outer sep=0pt},
%default radius will be 1pt. 
cross/.default={1pt}}

 % for braces
\usetikzlibrary{decorations.pathreplacing}
% for hashing area
\usetikzlibrary{patterns}
% tableaux var, signe
% source https://www.sqlpac.com/fr/documents/latex-package-tkz-tab-tikz-tableaux-de-signes-et-de-variations-de-fonctions.html
\usepackage{tkz-tab}

\tikzset{
	every axis/.style = {clip=true, grid style = {opacity=.5}}
}

\pagestyle{fancy}
\fancyhead[L]{Université de Mayotte}
\fancyhead[C]{\bf Série 1}
\fancyhead[R]{\today}

%%%%%%%%%%%%%%%%%%%%%%%%%%%%%%
% SELF MADE COLORS
%%%%%%%%%%%%%%%%%%%%%%%%%%%%%%

%!TEX encoding = UTF8
%%%%%%%%%%%%%%%%%%%%%%%%%%%%%%
% SELF MADE COLORS
%%%%%%%%%%%%%%%%%%%%%%%%%%%%%%


\definecolor{myg}{RGB}{56, 140, 70}
\definecolor{myb}{RGB}{45, 111, 177}
\definecolor{myr}{RGB}{199, 68, 64}
\definecolor{mytheorembg}{HTML}{F2F2F9}
\definecolor{mytheoremfr}{HTML}{00007B}
\definecolor{mylenmabg}{HTML}{FFFAF8}
\definecolor{mylenmafr}{HTML}{983b0f}
\definecolor{mypropbg}{HTML}{f2fbfc}
\definecolor{mypropfr}{HTML}{191971}
\definecolor{myexamplebg}{HTML}{F2FBF8}
\definecolor{myexamplefr}{HTML}{88D6D1}
\definecolor{myexampleti}{HTML}{2A7F7F}
\definecolor{mydefinitbg}{HTML}{E5E5FF}
\definecolor{mydefinitfr}{HTML}{3F3FA3}
\definecolor{notesgreen}{RGB}{0,162,0}
\definecolor{myp}{RGB}{197, 92, 212}
\definecolor{mygr}{HTML}{2C3338}
\definecolor{myred}{RGB}{127,0,0}
\definecolor{myyellow}{RGB}{169,121,69}
\definecolor{myexercisebg}{HTML}{F2FBF8}
\definecolor{myexercisefg}{HTML}{88D6D1}
\definecolor{doc}{RGB}{0,60,110}

% manim colors because they're beautiful
% https://docs.manim.community/en/stable/reference/manim.utils.color.manim_colors.html

\definecolor{BLACK}{HTML}{000000}\definecolor{BLUE}{HTML}{58C4DD}\definecolor{BLUE_A}{HTML}{C7E9F1}\definecolor{BLUE_B}{HTML}{9CDCEB}\definecolor{BLUE_C}{HTML}{58C4DD}\definecolor{BLUE_D}{HTML}{29ABCA}\definecolor{BLUE_E}{HTML}{236B8E}\definecolor{DARKER_GRAY}{HTML}{222222}\definecolor{DARKER_GREY}{HTML}{222222}\definecolor{DARK_BLUE}{HTML}{236B8E}\definecolor{DARK_BROWN}{HTML}{8B4513}\definecolor{DARK_GRAY}{HTML}{444444}\definecolor{DARK_GREY}{HTML}{444444}\definecolor{GOLD}{HTML}{F0AC5F}\definecolor{GOLD_A}{HTML}{F7C797}\definecolor{GOLD_B}{HTML}{F9B775}\definecolor{GOLD_C}{HTML}{F0AC5F}\definecolor{GOLD_D}{HTML}{E1A158}\definecolor{GOLD_E}{HTML}{C78D46}\definecolor{GRAY}{HTML}{888888}\definecolor{GRAY_A}{HTML}{DDDDDD}\definecolor{GRAY_B}{HTML}{BBBBBB}\definecolor{GRAY_BROWN}{HTML}{736357}\definecolor{GRAY_C}{HTML}{888888}\definecolor{GRAY_D}{HTML}{444444}\definecolor{GRAY_E}{HTML}{222222}\definecolor{GREEN}{HTML}{83C167}\definecolor{GREEN_A}{HTML}{C9E2AE}\definecolor{GREEN_B}{HTML}{A6CF8C}\definecolor{GREEN_C}{HTML}{83C167}\definecolor{GREEN_D}{HTML}{77B05D}\definecolor{GREEN_E}{HTML}{699C52}\definecolor{GREY}{HTML}{888888}\definecolor{GREY_A}{HTML}{DDDDDD}\definecolor{GREY_B}{HTML}{BBBBBB}\definecolor{GREY_BROWN}{HTML}{736357}\definecolor{GREY_C}{HTML}{888888}\definecolor{GREY_D}{HTML}{444444}\definecolor{GREY_E}{HTML}{222222}\definecolor{LIGHTER_GRAY}{HTML}{DDDDDD}\definecolor{LIGHTER_GREY}{HTML}{DDDDDD}\definecolor{LIGHT_BROWN}{HTML}{CD853F}\definecolor{LIGHT_GRAY}{HTML}{BBBBBB}\definecolor{LIGHT_GREY}{HTML}{BBBBBB}\definecolor{LIGHT_PINK}{HTML}{DC75CD}\definecolor{LOGO_BLACK}{HTML}{343434}\definecolor{LOGO_BLUE}{HTML}{525893}\definecolor{LOGO_GREEN}{HTML}{87C2A5}\definecolor{LOGO_RED}{HTML}{E07A5F}\definecolor{LOGO_WHITE}{HTML}{ECE7E2}\definecolor{MAROON}{HTML}{C55F73}\definecolor{MAROON_A}{HTML}{ECABC1}\definecolor{MAROON_B}{HTML}{EC92AB}\definecolor{MAROON_C}{HTML}{C55F73}\definecolor{MAROON_D}{HTML}{A24D61}\definecolor{MAROON_E}{HTML}{94424F}\definecolor{ORANGE}{HTML}{FF862F}\definecolor{PINK}{HTML}{D147BD}\definecolor{PURE_BLUE}{HTML}{0000FF}\definecolor{PURE_GREEN}{HTML}{00FF00}\definecolor{PURE_RED}{HTML}{FF0000}\definecolor{PURPLE}{HTML}{9A72AC}\definecolor{PURPLE_A}{HTML}{CAA3E8}\definecolor{PURPLE_B}{HTML}{B189C6}\definecolor{PURPLE_C}{HTML}{9A72AC}\definecolor{PURPLE_D}{HTML}{715582}\definecolor{PURPLE_E}{HTML}{644172}\definecolor{RED}{HTML}{FC6255}\definecolor{RED_A}{HTML}{F7A1A3}\definecolor{RED_B}{HTML}{FF8080}\definecolor{RED_C}{HTML}{FC6255}\definecolor{RED_D}{HTML}{E65A4C}\definecolor{RED_E}{HTML}{CF5044}\definecolor{TEAL}{HTML}{5CD0B3}\definecolor{TEAL_A}{HTML}{ACEAD7}\definecolor{TEAL_B}{HTML}{76DDC0}\definecolor{TEAL_C}{HTML}{5CD0B3}\definecolor{TEAL_D}{HTML}{55C1A7}\definecolor{TEAL_E}{HTML}{49A88F}\definecolor{WHITE}{HTML}{FFFFFF}\definecolor{YELLOW}{HTML}{FFFF00}\definecolor{YELLOW_A}{HTML}{FFF1B6}\definecolor{YELLOW_B}{HTML}{FFEA94}\definecolor{YELLOW_C}{HTML}{FFFF00}\definecolor{YELLOW_D}{HTML}{F4D345}\definecolor{YELLOW_E}{HTML}{E8C11C}

%%%%%%%%%%%%%%%%%%%%%%%%%%%%
% ENVIRONMENTS
%%%%%%%%%%%%%%%%%%%%%%%%%%%%

\theoremstyle{definition}

\newtheorem{theorem}{Théorème}
\newtheorem{corollaire}[theorem]{Corollaire}
\newtheorem{lemme}[theorem]{Lemme}
\newtheorem{proposition}[theorem]{Proposition}
\newtheorem{exercice}[theorem]{Exercice}
\newtheorem{exemple}[theorem]{Exemple}
\newtheorem{definition}[theorem]{Définition}
\newtheorem*{question}{Question}
\newtheorem*{preuve}{Preuve}
\newtheorem*{remarque}{Remarque}
\newtheorem*{strategie}{Stratégie}
\newtheorem*{methode}{Méthode}
\newtheorem*{notation}{Notation}
\newtheorem*{nomenclature}{Nomenclature}
\newtheorem{axiome}[theorem]{Axiome}
\newtheorem*{heuristique}{Heuristique}
\newtheorem*{motivation}{Motivation}

\newtheorem*{definition*}{Définition}
\newtheorem*{lemme*}{Lemme}
\newtheorem*{proposition*}{Proposition}
\newtheorem*{theorem*}{Théorème}
\newtheorem*{corollaire*}{Corollaire}

%%%%%%%%%%%%%%%%%%%%%%%%%%%%
% LETTERFONTS
%%%%%%%%%%%%%%%%%%%%%%%%%%%%

%!TEX encoding = UTF8

%fonts
\usepackage{libertinus,libertinust1math}
\usepackage[T1]{fontenc}

% for bold face math mode in section titles
\AddToHook{bfseries}{\boldmath}

% for calligraphic C, D, P (important to import this after the font)
\usepackage{calrsfs}
\newcommand{\D}{\mathcal{D}}
\newcommand{\Ccal}{\mathcal{C}}
\renewcommand{\P}{\mathcal{P}}

% Schwartz
\renewcommand{\S}{\mathcal{S}} % \S est le signe paragraphe normalement

% corps
\newcommand{\C}{\mathbb{C}}
\newcommand{\R}{\mathbb{R}}
\newcommand{\Rnn}{\mathbb{R}^{2n}}
\newcommand{\Z}{\mathbb{Z}}
\newcommand{\N}{\mathbb{N}}
\newcommand{\Q}{\mathbb{Q}}
\newcommand{\E}{\mathbb{E}}
\newcommand{\DD}{\mathbb{D}}
\newcommand{\K}{\mathbb{K}}
\newcommand{\F}{\mathbb{F}}
\newcommand{\U}{\mathbb{U}}
\renewcommand{\H}{\mathbb{H}}

% order notations
\DeclareRobustCommand{\O}{%
  \text{\usefont{OMS}{cmsy}{m}{n}O}%
}

% japanese bracket
\newcommand{\japb}[1]{\langle #1 \rangle}

% arrows over partial derivatives
\newcommand{\lp}{\overleftarrow{\partial}}
\newcommand{\rp}{\overrightarrow{\partial}}

% quantization
\newcommand{\h}{\hbar}
\newcommand{\Opht}{\textrm{Op}_{\h}^{t}}
\newcommand{\Op}[2][\hbar]{\textrm{Op}_{#1}^{#2}}

% omega functions
\newcommand{\omegap}[2][\rho_0]{\omega(\partial_{#1},\partial_{#2})}
\newcommand{\omegar}[2][\rho_0]{\omega(#1,#2)}

% space before semicolon
\mathcode`\;="303B

% for \Lightning
\usepackage{marvosym}

% for \warning
% 66 or 49 idk ; depends on the computer for some reason
\newcommand{\warning}{{\fontencoding{U}\fontfamily{futs}\selectfont\char 49\relax}}

% Q(\sqrt(d)) field
\newcommand{\Qsqrt}[1]{\Q\bigl(\mspace{-3mu}\sqrt{#1}\bigr)}

% closed off square root from https://en.wikibooks.org/wiki/LaTeX/Mathematics#Roots
% simple version: does not allow cbrt{x} = sqrt[3]{x}
% New definition of square root:
% it renames \sqrt as \oldsqrt
\let\oldsqrt\sqrt
% it defines the new \sqrt in terms of the old one
\def\sqrt{\mathpalette\DHLhksqrt}
\def\DHLhksqrt#1#2{%
\setbox0=\hbox{$#1\oldsqrt{#2\,}$}\dimen0=\ht0
\advance\dimen0-0.2\ht0
\setbox2=\hbox{\vrule height\ht0 depth -\dimen0}%
{\box0\lower0.4pt\box2}}



%%%%%%%%%%%%%%%%%%%%%%%%%%%%
% MACROS
%%%%%%%%%%%%%%%%%%%%%%%%%%%%

%!TEX encoding = UTF8


%%%%%%%%%%%%%%%%%%%%%%%%%%%%
% THEOREMS
%%%%%%%%%%%%%%%%%%%%%%%%%%%%

\newcommand{\thm}[3]{\begin{theorem}[#1]\label{#3}#2\end{theorem}}
\newcommand{\cor}[3]{\begin{corollaire}[#1]\label{#3}#2\end{corollaire}}
\newcommand{\lem}[3]{\begin{lemme}[#1]\label{#3}#2\end{lemme}}
\newcommand{\prop}[3]{\begin{proposition}[#1]\label{#3}#2\end{proposition}}
\newcommand{\ex}[3]{\begin{exemple}[#1]\label{#3}#2\end{exemple}}
\newcommand{\dfn}[3]{\begin{definition}[#1]\label{#3}#2\end{definition}}
\newcommand{\qs}[2]{\begin{question}[#1]#2\end{question}}

\renewcommand\qedsymbol{$\blacksquare$}
\newcommand{\pf}[2]{\begin{preuve}[#1]#2\hfill\qedsymbol\end{preuve}}

\newcommand{\nt}[1]{\begin{remarque}#1\end{remarque}}
\newcommand{\str}[1]{\begin{strategie}#1\end{strategie}}
\newcommand{\mth}[1]{\begin{methode}#1\end{methode}}
\newcommand{\ax}[3]{\begin{axiome}[#1]\label{#3}#2\end{axiome}}
\newcommand{\notations}[1]{\begin{notation}#1 \end{notation}}
\newcommand{\nomen}[1]{\begin{nomenclature}#1 \end{nomenclature}}
\newcommand{\heur}[1]{\begin{heuristique}#1\end{heuristique}}
\newcommand{\motiv}[1]{\begin{motivation}#1\end{motivation}}


%%%%%%%%%%%%%%%%%%%%%%%%%%%%
% EXERCISES
%%%%%%%%%%%%%%%%%%%%%%%%%%%%

\newcommand{\exe}[4]{
	\begin{Exercise}[#1, label=#3]
		%\marginpar{\mbox{\scriptsize(solution p.\pageref{\ExerciseLabel-Answer})}}
		#2
	\end{Exercise}
	\begin{Answer}[ref=#3]
		#4
	\end{Answer}
}

% version 1
\newcommand{\exemulticols}[5]{
	\begin{multicols}{2}
	\begin{Exercise}[#1, label=#4]
		#2
	\end{Exercise}
	\columnbreak
		#3
	\end{multicols}
	\begin{Answer}[ref=#4]
		#5
	\end{Answer}
}

% version 2
%\newcommand{\exemulticols}[5]{
%	\begin{Exercise}[#1, label=#4]
%	\begin{multicols}{2}
%		#2
%	\columnbreak
%		#3
%	\end{multicols}
%	\end{Exercise}
%	\begin{Answer}[ref=#4]
%		#5
%	\end{Answer}
%}

\newcommand{\exeTF}[1]{}

%%%%%%%%%%%%%%%%%%%%%%%%%%%%
% OPERATORS
%%%%%%%%%%%%%%%%%%%%%%%%%%%%

% deliminators
\DeclarePairedDelimiter{\abs}{\lvert}{\rvert}
\DeclarePairedDelimiter{\norm}{\lVert}{\rVert}
\DeclarePairedDelimiter{\scal}{\langle}{\rangle}

\DeclarePairedDelimiter{\ceil}{\lceil}{\rceil}
\DeclarePairedDelimiter{\floor}{\lfloor}{\rfloor}
\DeclarePairedDelimiter{\round}{\lfloor}{\rceil}

%  - others
\DeclareMathOperator{\Lap}{\mathcal{L}}
\DeclareMathOperator{\Var}{Var} % variance
\DeclareMathOperator{\Cov}{Cov} % covariance
\DeclareMathOperator{\proj}{proj} % projection
\DeclareMathOperator{\sym}{sym} % symétrique

\DeclareMathOperator{\RE}{Re} % partie réelle
\DeclareMathOperator{\IM}{Im} % partie imaginaire

\DeclareMathOperator{\pgcd}{pgcd}
\DeclareMathOperator{\ppcm}{ppcm}

\DeclareMathOperator{\ev}{ev}

\DeclareMathOperator{\tr}{tr}
\DeclareMathOperator{\sgn}{sgn}

\DeclareMathOperator{\Frac}{Frac}

% Lois
\DeclareMathOperator{\Bern}{\text{Bern}}
\DeclareMathOperator{\Binom}{\text{Binom}}

% I prefer the slanted \leq
\let\oldleq\leq % save them in case they're every wanted
\let\oldgeq\geq
\renewcommand{\leq}{\leqslant}
\renewcommand{\geq}{\geqslant}


%%%%%%%%%%%%%%%%%%%%%%%%%%%%
% commands
%%%%%%%%%%%%%%%%%%%%%%%%%%%%

% tel que
\newcommand{\tqs}{\text{ tels que }}
\newcommand{\tq}{\text{ tq. }}
\newcommand{\et}{\text{ et }}
\newcommand{\ou}{\text{ ou }}
\newcommand{\pourtout}{\text{ pour tout }}
\newcommand{\sct}{\text{ sachant }}
\newcommand{\id}{\text{id}}
\newcommand{\apriori}{\emph{a priori}}


% ensemble avec bigl et bigr
\newcommand{\bigset}[1]{\bigl\{ #1 \bigr\}}
\newcommand{\Bigset}[1]{\Bigl\{ #1 \Bigr\}}
\newcommand{\bigpar}[1]{\bigl( #1 \bigr)}
\newcommand{\Bigpar}[1]{\Bigl( #1 \Bigr)}

% PLUS INFTY AND MINUS INFTY WITH NO SPACE
\newcommand{\pinfty}{{+}\infty}
\newcommand{\minfty}{{-}\infty}

% vecteur flèche
\renewcommand{\vec}[1]{\overrightarrow{#1}}

% vecteur pmatrix
\newcommand{\pvec}[2]{\begin{pmatrix} #1 \\ #2 \end{pmatrix}}

% vecteur norme
%\newcommand{\norm}[1]{\left\Vert #1 \right\Vert}

% point plan
\newcommand{\point}[3]{
	#1\left(#2 ; #3 \right)
}

% \smash avant \underline pour coller la ligne au mot
% attention : ça empêche les line breaks donc c'est pas super...
\let\oldunderline\underline
\renewcommand{\underline}[1]{\oldunderline{\smash{#1}}}

% emph + index
% for \uline which allows underline with line breaks
\usepackage[normalem]{ulem}
% i add underline just for now
\newcommand{\emphindex}[1]{\emph{\uline{#1}}\index{#1}}

% tableau croisé
\newcommand{\tableaucroise}[4]{
\begin{tabular}{|c|c|c|c|}
	\cline{2-4}
	\multicolumn{1}{c|}{} & #1 \\ \hline
	#2 \\ \hline
	#3  \\ \hline
	#4  \\ \hline
\end{tabular}
}

% emdash in mathmode
\newcommand{\dash}{\text{---}}

%%%%%%%%%%%%%%%%%%%%%%%%%%%%
% ENVIRONMENTS
%%%%%%%%%%%%%%%%%%%%%%%%%%%%

% two columns toc
\newcommand{\twocolstableofcontents}{
	\renewcommand*\contentsname{}
	\begin{center}  
		{\Large \bfseries TABLE DES MATIÈRES \bigskip}    
	\end{center}
	\hrule
	\begin{multicols}{2}\setlength{\columnseprule}{1pt}
		\titleformat{\chapter}{}{}{0pt}{} 
		\titlespacing*{\chapter}{0pt}
		{*7} %  % vertical space before the title <<<<<<<<<<
		{*-15}  %  idem after title (in ex units + glue) <<<<<<<<<<<
		\tableofcontents
	\end{multicols}
	\hrule
}

% python minted
\newcommand{\python}[1]{
\inputminted[
		linenos,
		gobble=0,
		breaklines=true, % otherwise it breaks for no apparent reason?
		breakafter=,,
		fontsize=\normalsize,
		numbersep=8pt,
		tabsize=4, % tab ident = 4 spaces
		fontfamily=courier, %important pour les signes <, >
]{python}{python/#1.py}
}



%%%%%%%%%%%%%%%%%%%%%%%%%%%%
% PACKAGE-SPECIFIC ENVIRONMENTS
%%%%%%%%%%%%%%%%%%%%%%%%%%%%

% global settings for the first level in enumerate environments
\setlist[enumerate, 1]{%
	align = left, labelwidth = \parindent%, labelsep = 0pt, leftmargin = \parindent
}


%%%%%%%%%%%%%%%%%%%%%%%%%%%%%%
% EXERCISES
%%%%%%%%%%%%%%%%%%%%%%%%%%%%%%

\usepackage{ifthen}
\renewcommand{\ExerciseHeader}{
	\tikz[baseline=(R.base)]\node[draw,rectangle, thick, inner sep=2pt](R) {%
		\textbf{%
			\refAnswer{\ExerciseLabel}.%
		}%
	};\!
	%
	\ifnum\ExerciseDifficulty=0
	\else
		(\theExerciseDifficulty)
	\fi
}
\renewcommand{\DifficultyMarker}{$\star$}
\renewcommand{\AnswerHeader}{
	\tikz[baseline=(R.base)]\node[draw,rectangle, thick, inner sep=2pt](R) {\textbf{\theExercise.}};\!
}
\renewcommand{\exe}[4]{
	\begin{Exercise}[#1, label=#3]
		\if\relax\detokenize\expandafter{\ExerciseTitle}\relax
		%\marginpar{[Bonus]}
		\else
		\marginpar{\mbox{[\ExerciseTitle]}}
		\fi
		#2
	\end{Exercise}
	\begin{Answer}[ref=#3]
		#4
	\end{Answer}
}
\renewcommand{\exemulticols}[5]{
	\begin{Exercise}[#1, label=#4]
		\if\relax\detokenize\expandafter{\ExerciseTitle}\relax
		%\marginnote{[Bonus]}
		\else
		\marginnote{\mbox{[\ExerciseTitle]\qquad}}
		\fi
	\begin{multicols}{2}
		#2
	\columnbreak
		#3
	\end{multicols}
	\end{Exercise}
	\begin{Answer}[ref=#4]
		#5
	\end{Answer}
}





\SetDate[30/09/2026]

\begin{document}
\fancyhead[C]{\textbf{MG105 --- Fonctions de référence}}


	%\null\vspace{-40pt}


%\setlength{\columnsep}{2.1cm}
\setlength{\columnseprule}{0.4pt}
\newcommand{\scale}{.99}

\begin{multicols}{2}
	{
	\subsection*{Fonction affine croissante}
	\noindent
	Définition algébrique : $f(x) = $
	\\
	Domaine de définition : $\D_f = $
	
	\begin{center}
	\begin{tikzpicture}[scale=\scale]
		\begin{axis}[xmin = -4.1, xmax=4.1, ymin=-4.1, ymax=4.1, axis x line=middle, axis y line=middle, axis line style=->, grid=both, ticks=none,
	    	]
		\end{axis}
	\end{tikzpicture}	
	
	\vspace{10pt}
	
	\begin{tikzpicture}
		\tkzTabInit
		 [espcl=5]
	       		{$x$ / 1 , Variation de $f(x)$ / 2, Signe de $f(x)$ / 1}
	       		{$\minfty$,$\pinfty$}
	\end{tikzpicture}
	\end{center}
	}
	
	
	{
	\subsection*{Fonction carré}	
	\noindent
	Définition algébrique : $f(x) = $
	\\
	Domaine de définition : $\D_f = $
	
	\begin{center}
	\begin{tikzpicture}[scale=\scale]
		\begin{axis}[xmin = -5.1, xmax=5.1, ymin=-1.1, ymax=30.5, axis x line=middle, axis y line=middle, axis line style=->, grid=both,
		%ytick={-10,-8,...,10},
		extra x ticks = {0},
	    	]
		\end{axis}
	\end{tikzpicture}	
	
	\vspace{10pt}
	
	\begin{tikzpicture}
		\tkzTabInit
		 [espcl=5]
	       		{$x$ / 1 , Variation de $f(x)$ / 2, Signe de $f(x)$ / 1}
	       		{$\minfty$,$\pinfty$}
	\end{tikzpicture}
	\end{center}
	}
	
	{
	\subsection*{Fonction affine décroissante}	
	\noindent
	Définition algébrique : $f(x) = $
	\\
	Domaine de définition : $\D_f = $
	
	\begin{center}
	\begin{tikzpicture}[scale=\scale]
		\begin{axis}[xmin = -4.1, xmax=4.1, ymin=-4.1, ymax=4.1, axis x line=middle, axis y line=middle, axis line style=->, grid=both, ticks=none,
	    	]
		\end{axis}
	\end{tikzpicture}	
	
	\vspace{10pt}
	
	\begin{tikzpicture}
		\tkzTabInit
		 [espcl=5]
	       		{$x$ / 1 , Variation de $f(x)$ / 2, Signe de $f(x)$ / 1}
	       		{$\minfty$,$\pinfty$}
	\end{tikzpicture}
	\end{center}
	}
	
	
	{
	\subsection*{Fonction racine carrée}	
	\noindent
	Définition algébrique : $f(x) = $
	\\
	Domaine de définition : $\D_f = $
	
	\begin{center}
	\begin{tikzpicture}[scale=\scale]
		\begin{axis}[xmin = -.1, xmax=10.1, ymin=-.1, ymax=4.1, axis x line=middle, axis y line=middle, axis line style=->, grid=both,
		ytick={-10,-8,...,10},
		extra x ticks = {0},
		extra y ticks = {0},
	    	]
		\end{axis}
	\end{tikzpicture}	
	
	\vspace{10pt}
	
	\begin{tikzpicture}
		\tkzTabInit
		 [espcl=5]
	       		{$x$ / 1 , Variation de $f(x)$ / 2, Signe de $f(x)$ / 1}
	       		{$0$ ,$\pinfty$}
	\end{tikzpicture}
	\end{center}
	}


\end{multicols}

\newpage


\begin{multicols}{2}

	
	
	{
	\subsection*{Fonction inverse}	
	\noindent
	Définition algébrique : $f(x) = $
	\\
	Domaine de définition : $\D_f = $

	\begin{center}
	\begin{tikzpicture}[scale=\scale]
		\begin{axis}[xmin = -3.1, xmax=3.1, ymin=-12.3, ymax=12.3, axis x line=middle, axis y line=middle, axis line style=->, grid=both,
		ytick={-12,-8,...,10},
	    	]
		\end{axis}
	\end{tikzpicture}	
	
	\vspace{10pt}
	
	\begin{tikzpicture}
		\tkzTabInit
		 [espcl=5]
	       		{$x$ / 1 , Variation de $f(x)$ / 2, Signe de $f(x)$ / 1}
	       		{,}%{$\minfty$,$\pinfty$}
	\end{tikzpicture}
	\end{center}
	}
	
	
	{
	\subsection*{Fonction exponentielle}
	\noindent
	Définition algébrique : $f(x) = $
	\\
	Domaine de définition : $\D_f = $

	\begin{center}
	\begin{tikzpicture}[scale=\scale]
		\begin{axis}[xmin = -3.1, xmax=3.1, ymin=-0.3, ymax=12.3, axis x line=middle, axis y line=middle, axis line style=->, grid=both,
		%ytick={-12,-8,...,10},
		extra x ticks = {0},
	    	]
		\end{axis}
	\end{tikzpicture}	
	
	\vspace{10pt}
	
	\begin{tikzpicture}
		\tkzTabInit
		 [espcl=5]
	       		{$x$ / 1 , Variation de $f(x)$ / 2, Signe de $f(x)$ / 1}
	       		{,}%{$\minfty$,$\pinfty$}
	\end{tikzpicture}
	\end{center}
	}
	
	
	{
	\subsection*{Fonction cube}	
	\noindent
	Définition algébrique : $f(x) = $
	\\
	Domaine de définition : $\D_f = $
	
	\begin{center}
	\begin{tikzpicture}[scale=\scale]
		\begin{axis}[xmin = -3.1, xmax=3.1, ymin=-30.5, ymax=30.5, axis x line=middle, axis y line=middle, axis line style=->, grid=both,
		ytick={-30,-20,...,30},
	    	]
		\end{axis}
	\end{tikzpicture}	
	
	\vspace{10pt}
	
	\begin{tikzpicture}
		\tkzTabInit
		 [espcl=5]
	       		{$x$ / 1 , Variation de $f(x)$ / 2, Signe de $f(x)$ / 1}
	       		{,}%{$\minfty$,$\pinfty$}
	\end{tikzpicture}
	\end{center}
	}
	
	
	{
	\subsection*{Fonction logarithme népérien}	
	\noindent
	Définition algébrique : $f(x) = $
	\\
	Domaine de définition : $\D_f = $

	\begin{center}
	\begin{tikzpicture}[scale=\scale]
		\begin{axis}[xmin = -.1, xmax=20.1, ymin=-9.3, ymax=4.3, axis x line=middle, axis y line=middle, axis line style=->, grid=both,
		%ytick={-12,-8,...,10},
		extra y ticks = {0},
	    	]
		\end{axis}
	\end{tikzpicture}	
	
	\vspace{10pt}
	
	\begin{tikzpicture}
		\tkzTabInit
		 [espcl=5]
	       		{$x$ / 1 , Variation de $f(x)$ / 2, Signe de $f(x)$ / 1}
	       		{,}%{$0$,$\pinfty$}
	\end{tikzpicture}
	\end{center}
	}

\end{multicols}



\newpage


\begin{multicols}{2}

	
	
	{
	\subsection*{Fonction valeur absolue}	
	\noindent
	Définition algébrique : $f(x) = $
	\\
	Domaine de définition : $\D_f = $

	\begin{center}
	\begin{tikzpicture}[scale=\scale]
		\begin{axis}[xmin = -3.1, xmax=3.1, ymin=-.3, ymax=12.3, axis x line=middle, axis y line=middle, axis line style=->, grid=both,
		%ytick={-12,-8,...,10},
		extra x ticks = {0},
	    	]
		\end{axis}
	\end{tikzpicture}	
	
	\vspace{10pt}
	
	\begin{tikzpicture}
		\tkzTabInit
		 [espcl=5]
	       		{$x$ / 1 , Variation de $f(x)$ / 2, Signe de $f(x)$ / 1}
	       		{,}%{$\minfty$,$\pinfty$}
	\end{tikzpicture}
	\end{center}
	}
	
	
	{
	\subsection*{Fonction cosinus}
	\noindent
	Définition algébrique : $f(x) = $
	\\
	Domaine de définition : $\D_f = $

	\begin{center}
	\begin{tikzpicture}[scale=\scale]
		\begin{axis}[xmin = -9.5, xmax=9.5, ymin=-1.1, ymax=1.1, axis x line=middle, axis y line=middle, axis line style=->, grid=both,
	%	ytick={-30,-20,...,30},
		xtick={-3*3.14159, -2*3.14159, -3.14159, 0, 3.14159, 2*3.14159, 3*3.14159},
		xticklabels={$-3\pi$, $-2\pi$, $-\pi$, $0$, $\pi$, $2\pi$, $3\pi$}
	    	]
		\end{axis}
	\end{tikzpicture}	
	
	\vspace{10pt}
	
	\begin{tikzpicture}
		\tkzTabInit
		 [espcl=5]
	       		{$x$ / 1 , Variation de $f(x)$ / 2, Signe de $f(x)$ / 1}
	       		{,}%{$\minfty$,$\pinfty$}
	\end{tikzpicture}
	\end{center}
	}
	
	
	{
	\subsection*{Fonction sinus}	
	\noindent
	Définition algébrique : $f(x) = $
	\\
	Domaine de définition : $\D_f = $
	
	\begin{center}
	\begin{tikzpicture}[scale=\scale]
		\begin{axis}[xmin = -9.5, xmax=9.5, ymin=-1.1, ymax=1.1, axis x line=middle, axis y line=middle, axis line style=->, grid=both,
	%	ytick={-30,-20,...,30},
		xtick={-3*3.14159, -2*3.14159, -3.14159, 0, 3.14159, 2*3.14159, 3*3.14159},
		xticklabels={$-3\pi$, $-2\pi$, $-\pi$, $0$, $\pi$, $2\pi$, $3\pi$}
	    	]
		\end{axis}
	\end{tikzpicture}	
	
	\vspace{10pt}
	
	\begin{tikzpicture}
		\tkzTabInit
		 [espcl=5]
	       		{$x$ / 1 , Variation de $f(x)$ / 2, Signe de $f(x)$ / 1}
	       		{,}%{$\minfty$,$\pinfty$}
	\end{tikzpicture}
	\end{center}
	}
	
	
	{
	\subsection*{Fonction tangente}	
	\noindent
	Définition algébrique : $f(x) = $
	\\
	Domaine de définition : $\D_f = $

	\begin{center}
	\begin{tikzpicture}[scale=\scale]
		\begin{axis}[xmin = -pi/2, xmax=pi/2, ymin=-10.1, ymax=10.1, axis x line=middle, axis y line=middle, axis line style=->, grid=both,
		xtick={-.5*3.14159, -.25*3.14159, 0, .25*3.14159, .5*3.14159},
		xticklabels={$-\frac{\pi}2$, $-\frac{\pi}4$, $0$, $\frac{\pi}4$, $\frac{\pi}2$},
		ytick distance = 2,
	    	]
		\end{axis}
	\end{tikzpicture}	
	
	\vspace{10pt}
	
	\begin{tikzpicture}
		\tkzTabInit
		 [espcl=5]
	       		{$x$ / 1 , Variation de $f(x)$ / 2, Signe de $f(x)$ / 1}
	       		{,}%{$\minfty$,$\pinfty$}
	\end{tikzpicture}
	\end{center}
	}

\end{multicols}


%\ExerciseSelect
%	[label={
%		exe:ref1,%
%		exe:ref2,%
%		exe:ref3,%
%		exe:ref4,%
%		exe:ref5,%
%		exe:ref6,%
%		exe:ref7,%
%		exe:ref8,%
%		%exe:ref9,%
%		exe:ref10,%
%		exe:ref11,%
%	}]
%%!TEX encoding = UTF8
%!TEX root = 

%%%%%%%%%%%%%%%%%%%%%%%%
%%%%%%%%%% SET 1 %%%%%%%%%%
%%%%%%%%%%%%%%%%%%%%%%%% 


\exe{}{
	Écrire les sommes et différences de rationnels sous forme de fraction irréductible.
	\begin{multicols}{3}
	\begin{enumerate}[itemsep=10pt]
		\item $1 - \frac12$
		\item $\frac78 + \frac{9}{40}$
		\item $\frac{3}{4} + \frac{2}{5}$
		\item $\frac{7}{6} - \frac{5}{8}$
		\item $\frac{9}{10} + \frac{2}{3} - \frac{1}{4}$
		\item $\frac{-5}{7} + \frac{8}{9}$
%		\item $\frac{4}{5} - \frac{3}{10} + \frac{-5}{6}$
%		\item $\frac{-1}{2} + \frac{7}{8}$
%		\item $\frac{11}{12} - \frac{5}{9} + \frac{7}{10}$
%		\item $\frac{2}{3} - \frac{5}{4}$
%		\item $\frac{-1}{3} + \frac{7}{9} - \frac{5}{6}$
%		\item $\frac{3}{8} - \frac{1}{2}$
	\end{enumerate}
	\end{multicols}
}{exe:rationnals}
{
	\begin{multicols}{2}
	\begin{enumerate}
		\item $1 - \frac{1}{2} = \frac{1}{2}$
		\item $\frac{7}{8} + \frac{9}{40} = \frac{11}{10}$
		\item $\frac{3}{4} + \frac{2}{5} = \frac{23}{20}$
		\item $\frac{7}{6} - \frac{5}{8} = \frac{13}{24}$
		\item $\frac{9}{10} + \frac{2}{3} - \frac{1}{4} = \frac{79}{60}$
		\item $\frac{-5}{7} + \frac{8}{9} = \frac{11}{63}$
%		\item $\frac{4}{5} - \frac{3}{10} + \frac{-5}{6} = \frac{-1}{3}$
%		\item $\frac{-1}{2} + \frac{7}{8} = \frac{3}{8}$
%		\item $\frac{11}{12} - \frac{5}{9} + \frac{7}{10} = \frac{191}{180}$
%		\item $\frac{2}{3} - \frac{5}{4} = \frac{-7}{12}$
%		\item $\frac{-1}{3} + \frac{7}{9} - \frac{5}{6} = \frac{-7}{18}$
%		\item $\frac{3}{8} - \frac{1}{2} = \frac{-1}{8}$
	\end{enumerate}
	\end{multicols}
}


\exe{}{
	Écrire les produits et divisions de rationnels sous forme de fraction irréductible.
	
	\begin{multicols}{3}
	\begin{enumerate}[itemsep=10pt]
		\item $\frac{\ \ -3 \ \ }{\frac54}$
		\item $\frac{63}{20} \times \frac{15}{14}$
		\item $\frac{3}{4} \times \frac{2}{5}$
		\item $\frac{\ \ \frac{7}{6}\ \ }{\frac{5}{8}}$
		\item $\frac{9}{10} \times \frac{2}{3} \times \frac{4}{1}$
		\item $\frac{-5}{7} \times \frac{8}{9}$
%		\item $\left(\frac{4}{5}\right)^{-1} \times \frac{3}{10}$
%		\item $\frac{-1}{2} \times \left(\frac{7}{8}\right)^{-1}$
%		\item $\frac{11}{12} \times \frac{5}{9} \times \left(\frac{7}{10}\right)^{-1}$
%		\item $\frac{2}{3} \times \frac{5}{4}$
%		\item $\frac{-1}{3} \times \left(\frac{7}{9}\right)^{-1}$
%		\item $\frac{3}{8} \times \frac{1}{2}$
	\end{enumerate}
	\end{multicols}
}{exe:rationnals2}
{
	\begin{multicols}{2}
	\begin{enumerate}[itemsep=10pt]
		\item $\frac{-3}{\frac{5}{4}} = \frac{-12}{5}$
		\item $\frac{63}{20} \times \frac{15}{14} = \frac{27}{8}$
		\item $\frac{3}{4} \times \frac{2}{5} = \frac{3}{10}$
		\item $\frac{\ \ \frac{7}{6}\ \ }{\frac{5}{8}} = \frac{28}{15}$
		\item $\frac{9}{10} \times \frac{2}{3} \times \frac{4}{1} = \frac{12}{5}$
		\item $\frac{-5}{7} \times \frac{8}{9} = \frac{-40}{63}$
%		\item $\left(\frac{4}{5}\right)^{-1} \times \frac{3}{10} = \frac{3}{8}$
%		\item $\frac{-1}{2} \times \left(\frac{7}{8}\right)^{-1} = \frac{-4}{7}$
%		\item $\frac{11}{12} \times \frac{5}{9} \times \left(\frac{7}{10}\right)^{-1} = \frac{275}{378}$
%		\item $\frac{2}{3} \times \frac{5}{4} = \frac{5}{6}$
%		\item $\frac{-1}{3} \times \left(\frac{7}{9}\right)^{-1} = \frac{-3}{7}$
%		\item $\frac{3}{8} \times \frac{1}{2} = \frac{3}{16}$
	\end{enumerate}
	\end{multicols}
}


\exe{}{
	Exprimer chaque valeur suivante sous forme de fraction irréductible.
	\begin{multicols}{4}
	\begin{enumerate}[label=\alph*), itemsep=10pt]
		\item $1-\frac{4}{44}$
		\item $\frac{15}{27}$
		\item $\frac{8}{3} + 6$
		\item $\frac{65}{5^2}$
		\item $\frac{6}{32}\times\frac{2}{3}$
		\item $\frac{38}{20}-3$
		\item $\left(\frac{10}{44}\right)^{-1}$
		\item $\frac{6}{18}\times2 + 1$
		\item $\left(\frac{66}{99}\right)^2$
		\item $\frac{23}{18} + \frac12$
		\item $\frac{34}{20}$
		\item $2 - \frac{9}{5}$
	\end{enumerate}
	\end{multicols}
}{exe:frac-irr}{
	\begin{multicols}{4}
	\begin{enumerate}[label=\alph*), itemsep=10pt]
		\item $\frac{10}{11}$
		\item $\frac{5}{9}$
		\item $\frac{26}{3}$
		\item $\frac{13}{5}$
		\item $\frac{1}{18}$
		\item $\frac{-11}{10}$
		\item $\frac{22}{5}$
		\item $\frac{5}{3}$
		\item $\frac{4}{9}$
		\item $\frac{16}{9}$
		\item $\frac{17}{10}$
		\item $\frac{1}{5}$
	\end{enumerate}
	\end{multicols}
}


\exe{difficulty=0}{
	En décembre pour les fêtes, M. Marchand dit avoir vendu les quatre cinquièmes de sa marchandise. En janvier, pendant les soldes, il a encore vendu les trois quarts de ce qu'il restait.
	\begin{enumerate}
		\item
		Quelle fraction de sa marchandise a-t-il vendu en tout?
		\item  
		La valeur totale de sa marchandise est de 262 000 euros. Quelle somme représente sa vente globale?
	\end{enumerate}

}{exe:Samoyeau6}{
	\begin{enumerate}
		\item
		D'abord, une fraction $\frac45$ de la marchandise est vendue.
		Il reste donc $1-\frac45 = \frac15$ de la marchandise.
		
		Ensuite, $\frac34$ de ce qu'il reste est vendu.
		L'expression « $\frac34$ de $\frac15$ » correspond au calcul $\frac34 \times \frac15 = \frac3{20}$.
		
		Finalement, $\frac45 + \frac3{20} = \frac{19}{20}$ de la marchandise est vendue en tout.
		\item  
		L'expression « $\frac{19}{20}$ de 262 000 » correspond au calcul $\frac{19}{20} \times 262~000$.
		Or, comme $\frac{19}{20} = 1 - \frac{1}{20}$, nous avons
			\[ \frac{19}{20}\times 262~000 = \left(1-\frac1{20}\right) \times 262~000 = 262~000 - \frac1{20}\times262~000. \]
		Il suffit donc de calculer $\frac1{20}$ de 262 000, ce qui n'est pas trop difficile car $\frac1{20} = \frac12 \times \frac1{10}$, il suffit donc de diviser par 2 puis par 10.
		Ainsi, $\frac1{20} \times 262~000 = 13~100$ et $\frac{19}{20}\times 262~000 = 262~000 - 13~100 = 248~900$.
	\end{enumerate}
}

\exe{difficulty=0}{
	Une usine française de voitures vend $\frac{1}{3}$ de ses voitures en Allemagne, $\frac{1}{4}$ de ce qui reste en Espagne et le reste en France. Quelle proportion de voitures de cette usine reste en France?
}{exe:Samoyeau7}{
	D'abord, $\frac13$ est vendu, et il reste $1-\frac13 = \frac23$.
	
	Ensuite, « $\frac14$ de $\frac23$ » correspond au calcul $\frac14 \times \frac23 = \frac16$, proportion de vendue en plus.
	
	Au totale, la proportion vendue est égale à $\frac13 + \frac16 = \frac12$, et donc la proportion invendue est $1-\frac12 = \frac12$.
}

\exe{difficulty=0}{
	Amaya et Mariza ont deux tablettes de chocolat, une au chocolat noir et une au chocolat au lait qui sont de même proportion. 
	Amaya a mangé $\frac{1}{4}$ des $\frac{2}{3}$ de la tablette au chocolat au lait. Mariza a mangé la moitié des $\frac{3}{4}$ de la tablette de chocolat noir. Laquelle a mangé le plus de chocolat? 
}{exe:Samoyeau8}{
	Amaya a mangé « $\frac14$ de $\frac23$ », ce qui correspond au calcul $\frac14\times\frac23 = \frac16$.
	
	Mariza a mangé « $\frac12$ de $\frac34$ », ce qui correspond au calcul $\frac12 \times \frac34 = \frac38$.
	
	Pour comparer deux fractions, il suffit de les mettre au même dénominateur et de comparer les numérateurs.
	Ici, en mettant sur $24$, on trouve
		\[ \frac16 = \frac4{24} < \frac{9}{24} = \frac38, \]
	car $4 < 9$.
	
	En conclusion, Mariza a mangé plus de chocolat car les tablettes sont supposées identiques au départ.
}

\exe{difficulty=0}{
	Lors d'une course, Pierre fait $\frac{1}{3}$ de la distance totale parcourue en vélo, $\frac{5}{12}$ de la distance totale en courant, $\frac{1}{6}$ de la distance totale à la nage et le reste en kayak. 
	\begin{enumerate}
		\item Compléter la phrase suivante : Pierre a parcouru $\dots\dots$ de la distance totale en kayak.
		\item Cette course fait 24 kilomètres. Quelle est la distance parcourue en kayak? 
	\end{enumerate} 

}{exe:Samoyeau9}{
	\begin{enumerate}
		\item  
		La somme des proportion fait $1$, donc le reste est donné par $1 - \frac13 - \frac5{12} - \frac16 = \frac1{12}$.
		
		Pierre a parcouru $\frac1{12}$ de la distance totale en kayak.
		\item 
		L'expression « $\frac1{12}$ de $24$ » correspond au calcul $\frac1{12} \times 24 = 2$.
		La distance est donc égale à 2km.
	\end{enumerate} 
}

\exe{difficulty=0}{
	Considérons le nombre zêta : 
	$ 
	\zeta=\frac{ \frac{2}{3} \times \frac{6}{12} + \frac{-3}{7} - 4 }{ \frac{\frac{1}{2}-3}{8\times \frac{-1}{2}} + \frac{2}{5}-\frac{1}{4}}
	$.
	L'exprimer sous forme de fraction irréductible.

}{exe:Samoyeau10}{
	WolframAlpha arrive ici à la rescousse et nous donne $\zeta = 	-\frac{3400}{651}$, résultat qu'on croira aveuglément.
}

\exe{difficulty=0}{
	Classer les fractions suivantes par ordre croissant : 
	\begin{align*}
		A=\frac{1}{2}, && B=\frac{1}{6}, && C=\frac{5}{12}, &&D=\frac{3}{4}, &&E=\frac{2}{3}.
	\end{align*}
}{exe:Samoyeau2}{
	Pour comparer deux fractions, il suffit de les mettre au même dénominateur et de comparer les numérateurs.
	Ici, en mettant sur $24$, on trouve
	\begin{align*}
		A=\frac{12}{24}, && B=\frac{4}{24}, && C=\frac{10}{24}, &&D=\frac{18}{24}, &&E=\frac{16}{24}.
	\end{align*}
	Il suit donc que $B < C < A < E < D$ par inspection des numérateurs.
}


\exe{}{
	Sans calculatrice, donner le développement décimal des fractions suivantes.
	\begin{multicols}{3}
	\begin{enumerate}[label=\roman*), leftmargin=60pt]
		\item $\frac1{10}$
		\item $\frac1{5}$
		\item $\frac3{10^{5}}$
		\item $\frac7{20}$
		\item $\frac{395}{50}$
		\item $\frac{11}{200}$
	\end{enumerate}
	\end{multicols}
}{exe:dev-decimaux}{
	\begin{multicols}{3}
	\begin{enumerate}[label=\roman*)]
		\item $\frac1{10} = 0,1.$
		\item $\frac1{5} = \frac2{10} = 0,2.$
		\item $\frac3{10^{5}} = 0,000~03.$
		\item $\frac7{20} = \frac{3,5}{10} = 0,35.$
		\item $\frac{395}{50} = \frac{790}{100} = 7,9.$
		\item $\frac{11}{200} = \frac{5,5}{100} = 0,055.$
	\end{enumerate}
	\end{multicols}
}


\exe{}{
	Sans calculatrice, donner le développement décimal des fractions suivantes.
	\begin{multicols}{3}
	\begin{enumerate}[label=\roman*), leftmargin=60pt]
		\item $\frac1{100}$
		\item $\frac3{5}$
		\item $\frac2{10^{1}}$
		\item $\frac7{2}$
		\item $\frac{7}{20}$
		\item $\frac{31}{200}$
	\end{enumerate}
	\end{multicols}
}{exe:dev-decimauxE}{
	\begin{multicols}{3}
	\begin{enumerate}[label=\roman*)]
		\item $\frac1{100} = 0,01$
		\item $\frac3{5} = 0,6$
		\item $\frac2{10^{1}} = 0,2$
		\item $\frac7{2} = 3,5$
		\item $\frac{7}{20} = 0,35$
		\item $\frac{31}{200} = 0,155$
	\end{enumerate}
	\end{multicols}
}


\exe{, difficulty=1}{
	Déterminer le développement décimal de $\frac1{11}$ et de $\frac{12}{22}$ en posant la division.
}{exe:111-dec}{
	Posons 
		\[ \frac{1}{11} = 0,n_1 n_2 n_3 n_4 \dots. \]
	Par multplication par 10 successives et étude des parties décimales, on obtient $n_1 = 0$, $n_2=9$, et $\frac1{11} = 0,n_3n_4n_5\dots$.
	D'où
		\[ \frac1{11} = 0,09~09~09~09~\cdots. \]
	Posons 
		\[ \frac{12}{22} = 0,n_1 n_2 n_3 n_4 \dots. \]
	Par multplication par 10 successives et étude des parties décimales, on obtient $n_1 = 5$, $n_2=4$, et $\frac{12}{22} = 0,n_3 n_4 n_5\dots$.
	D'où
		\[ \frac{12}{22} = 0,54~54~54~54~\cdots. \]
}



\exe{, difficulty=1}{
	Montrer que $\sqrt2 = 1 + \frac1{1+\sqrt2}$.
	En déduire que $\sqrt2 = 1 + \frac1{2 + \frac1{1+\sqrt2}}$ et que
		\[\sqrt2 = 1 + \cfrac1{2 + \cfrac1{2 + \cfrac1{2 + \cfrac1{1+\sqrt2}}}}. \]
	Exprimer le rationnel $1 + \frac1{2 + \frac1{2 + \frac12}}$ sous la forme $\frac{p}{q}$ avec $p$ et $q$ premiers entre eux et calculer $\left| \sqrt2 - \frac{p}q \right|$ à l'aide de la calculatrice.
	Faire idem avec $1 + \frac1{2 + \frac1{2 + \frac1{2 + \frac12}}}$.
}{exe:frac-continue-sqrt2}{
	En manipulant l'équation on se retrouve à vouloir montrer que $\sqrt2 - 1 = \frac1{\sqrt2+1}$, ce qui est équivalent à $(\sqrt2 - 1)(\sqrt2 + 1) = 1$, vrai par identité remarquable (ou par distributivité simple).
	
	Comme $\sqrt2 = 1 + \frac1{1+\sqrt2}$, on peut remplacer le $\sqrt2$ dans l'expression de droite par toute l'expression de droite.
	On obtient donc $\sqrt2 = 1 + \cfrac1{1 + 1 + \cfrac1{1+\sqrt2}} = 1 + \cfrac1{2 + \cfrac1{1+\sqrt2}}$ comme voulu.
	On obtient l'expression d'après en répétant le même raisonnement.
	
	Par simplifications simples, on trouve $\frac{p}{q} = \frac{17}{12}$.
	La calculatrice nous donne alors $\left| \sqrt2 - \frac{17}{12} \right| \approx 0,002$.
	
	Finalement, remarquons que 
	\begin{align*}
		1 + \cfrac1{2 + \cfrac1{2 + \cfrac1{2 + \cfrac12}}} &= 1 + \cfrac1{1+ 1 + \cfrac1{2 + \cfrac1{2 + \cfrac12}}}, \\
			&= 1 + \cfrac1{1+ \frac{17}{12}} = \frac{41}{29}.
	\end{align*}
	Le rationnel $\frac{41}{29}$ vérifie $\left| \sqrt2 - \frac{41}{29} \right| \approx 0,0004$, soit 5 fois plus proche de $\sqrt2$ que $\frac{17}{12}$ l'est.
}

\exe{difficulty=2}{
	En s'inspirant de l'exercice \ref{exe:frac-continue-sqrt2}, calculer sans calculatrice un rationnel s'approchant de $\sqrt5$ à $10^{-3}$ près.
}{exe:frac-continue-sqrt5}{
	Comme $2 < \sqrt5 < 3$, il est tentant d'écrire $\sqrt5 = 2 + \frac1{x}$ pour un certain $x\in\R$.
	Ce nombre $x$ vérifie donc
		\[ x = \frac1{\sqrt5 - 2} = \frac{\sqrt5+2}{(\sqrt5-2)(\sqrt5+2)} = \frac{\sqrt5+2}{5-4} =  \sqrt5 +2. \]
	Par conséquent
		\[ \sqrt5 = 2 + \frac1{2+\sqrt5} = 2 + \cfrac{1}{2+2+\cfrac{1}{2+\sqrt5}} = 2 + \cfrac{1}{4 + \cfrac1{2+\sqrt5}} = \dots \]
}

\exe{difficulty=1}{
	Vrai ou faux ? Démontrer que chaque proposition est vraie ou donner un contre-exemple.
	\begin{enumerate}
		\item
		La somme de deux nombres rationnels est un nombre rationnel.
		\item
		Le produit de deux nombres rationnels est un nombre rationnel.
		\item
		Un nombre rationnel à la puissance d'un nombre rationnel est un nombre rationnel.
		\item
		Un nombre rationnel à la puissance d'un nombre rationnel est un nombre irrationnel.
		\item
		La somme d'un nombre irrationnel et d'un nombre rationnel est un nombre irrationnel.
		\item
		La somme de deux nombres irrationnels est un nombre irrationnel.
		\item[7. $(\star\star)$]
		Un nombre irrationnel à la puissance d'un nombre irrationnel est irrationnel.
	\end{enumerate}
}{exe:VF-rationnels}{
	\begin{enumerate}
		\item
		Vrai.
		\item
		Vrai.
		\item
		Faux car $2^{\frac12} = \sqrt2$ qui n'est pas rattionnel.
		\item
		Faux car $2^2 = 4$ qui est rationnel.
		\item
		Vrai car si la somme était rationnelle, alors un nombre irrationnelle serait différence de deux rationnels et serait donc rationnel d'après le premier point.
		\item
		Faux car $\sqrt2 + (-\sqrt2) = 0$ est rationnel.
		\item[7. $(\star\star)$]
		Faux car soit $\sqrt2^{\sqrt2}$ est rationnel et donc un contre-exemple, soit $\left(\sqrt2^{\sqrt2}\right)^{\sqrt2} = 2$ est rationnel et donc un contre-exemple.
	\end{enumerate}
}


\exe{, difficulty=1}{
	En musique, la gamme classique contient douze notes avec un rapport de fréquence de 2 sur l'octave.
	La gamme est dite à tempérament égal si le rapport de fréquence entre deux notes consécutives est toujours le même.
	Notons ce rapport $r\in\R$.
	Justifier que $r^{12} = 2$ puis démontrer que $r$ est un nombre irrationnel : $r \not\in\Q$.
}{exe:temp-egal}{
	Le rapport de fréquence entre deux notes consécutives étant toujours $r$, on multiplie $r$ par lui-même 12 fois pour obtenir le rapport de fréquence entre deux notes de distance $12$, c'est-à-dire deux notes séparée par un octave.
	Ainsi $r^{12} = 2$ car l'octave correspond à un rapport de fréquence de 2.
	
	Supposons désormais que $r$ soit rationnel. Alors $r^{6}$ est aussi rationnel.
	Cependant, $(r^6)^2 = r^{12} = 2$, donc $r^{6} = \sqrt2$, qui n'est pas rationnel. \Large\Lightning
}



%%%%%%%%%%%%%%%%%%%%%%%%
%%%%%%%%%% SET 2 %%%%%%%%%%
%%%%%%%%%%%%%%%%%%%%%%%% 



\exe{}{
	Trouver le $x\in\Q$ rationnel vérifiant chaque équation suivante, et l'exprimer sous forme de fraction irréductible.

	\begin{multicols}{2}
	\begin{enumerate}[itemsep=10pt]
		\item $3x + 5 = 2x + 7$
		\item $4x - 6 = 3x + 8$
		\item $5x + 10 = 2x + 4$
		\item $-2x + 3 = 4x - 1$
		\item $7x - 5 = 5x + 9$
		\item $6x + 4 = 3x + 12$
		\item $-3x + 2 = -5x - 4$
		\item $8x - 9 = 6x + 3$
		\item $2x + 11 = -3x + 5$
		\item $-4x + 7 = -x + 10$
	\end{enumerate}
	\end{multicols}
}{exe:inter-affine}
{
	\begin{multicols}{2}
	\begin{enumerate}[itemsep=10pt]
		\item $3x + 5 = 2x + 7 \iff x = 2$
		\item $4x - 6 = 3x + 8 \iff x = 14$
		\item $5x + 10 = 2x + 4 \iff x = -2$
		\item $-2x + 3 = 4x - 1 \iff x = \frac{2}{3}$
		\item $7x - 5 = 5x + 9 \iff x = 7$
		\item $6x + 4 = 3x + 12 \iff x = \frac{8}{3}$
		\item $-3x + 2 = -5x - 4 \iff x = -3$
		\item $8x - 9 = 6x + 3 \iff x = 6$
		\item $2x + 11 = -3x + 5 \iff x = \frac{-6}{5}$
		\item $-4x + 7 = -x + 10 \iff x = -1$
	\end{enumerate}
	\end{multicols}
}



\exe{, difficulty=1}{
	Pour chaque paire d'équations, ajouter entre elles le symbole
		\begin{enumerate}[itemindent=1.5cm]
			\item[$\implies$] si la première équation implique la seconde ;
			\item[$\impliedby$] si la deuxième équation implique la première ;
			\item[$\iff$] si les deux équations sont équivalentes ;
			\item[$\centernot\iff$] si les deux équations sont indépendantes.
		\end{enumerate}
	Spécifier l'opération appliquée pour obtenir l'équation impliquée à partir de l'autre.
	
	
	\begin{center}
	\begin{tabular}{ccc|c|c}
		Équation 1 & Symbole & Équation 2 & Opération & \thead{Les équations sont-\\ elles équivalentes ?} \\ \hline
		$3x + 2y = 1$ &  & $6x  + 4y = 2$ & &  \\ \hline
		$-x + 2y = 1$ &  & $-x + 2y - 1 = 0$ & &  \\ \hline
		$3x + 2y = 1$ & & $-3x - 5y = -3$ & &  \\ \hline
		$x^2 = -3x$ &  & $x = -3$ & &  \\ \hline
		$-10x - y = 3$ &  & $0=0$ & &  \\ \hline
		$x-y= 3$ & & $-x + y = -3$ & &  \\ \hline
		$x -y = 0$ & & $x=y$ & &  \\ \hline
		$-x + 2y = 3$ & & $0=0$ & &  \\ \hline
		$2x + 3y + 4 = 2x + 3y + 5$ & & $1=0$ & &  \\ \hline
		$-x - y = 1$ & & $-x - y = 1$ & &  \\ \hline
		$0=0$ & & $23x - 10y = 6$ & &  \\ \hline
		$x + 2y = 3$ & & $x^2 + 2yx = 3x$ & &  \\ \hline
		$x-3y  =0$ & & $-2x+6y = 0$ & &  \\ \hline
		$-2x + 3y = -3$ & & $4x - 6y = -6$ & &  \\ \hline
		$x^2 = x$ & & $x = 1$ & &  \\ \hline
		$x = 3$ & & $x^2 = 9$ & &  \\ \hline
	\end{tabular}
	\end{center}
	
	{Remarque : par abus de notation et pour gagner du temps, l'absence de symbole signifie en général que les équations sont équivalentes, mais il faut toujours faire attention à ce que cela soit bien le cas !}
	
	\textbf{Attention !} si $f(x) = 0 \implies g(x) = 0$, alors une solution de $f(x) = 0$ est une solution de $g(x) = 0$.
	En revanche, une solution de $g(x) = 0$ peut ne pas être une solution de $f(x) = 0$.
	
}{exe:equations-equivalentes}{
	\begin{center}
	\begin{tabular}{ccc|c|c}
		Équation 1 & Symbole & Équation 2 & Opération & \thead{Les équations sont-\\ elles équivalentes ?} \\ \hline
		$3x + 2y = 1$ & {$\iff$} & $6x  + 4y = 2$ & $\times2$ ou $\times\frac12$ & {Oui} \\ \hline
		$-x + 2y = 1$ & {$\iff$} & $-x + 2y - 1 = 0$ & {$-1$ ou $+1$} & {Oui} \\ \hline
		$3x + 2y = 1$ & {$\centernot\iff$} & $-3x - 5y = -3$ & & {Non} \\ \hline
		$x^2 = -3x$ & {$\impliedby$} & $x = -3$ &$\times x$ & {Non} \\ \hline
		$-10x - y = 3$ & {$\implies$} & $0=0$ & {$\times0$} & {Non} \\ \hline
		$x-y= 3$ & {$\iff$} & $-x + y = -3$ & {$\times(-1)$} & {Oui} \\ \hline
		$x -y = 0$ &{$\iff$} & $x=y$ & {$+y$ ou $-y$} & {Oui} \\ \hline
		$-x + 2y = 3$ & {$\implies$} & $0=0$ & {$\times0$} & {Non} \\ \hline
		$2x + 3y + 4 = 2x + 3y + 5$ & {$\iff$} & $1=0$ & $+(-2x-3y)$ ou $-(-2x-3y)$ & {Oui} \\ \hline
		$-x - y = 1$ & {$\centernot\iff$} & $-x - y = 1$ & & {Non} \\ \hline
		$0=0$ & {$\impliedby$} & $23x - 10y = 6$ & $\times0$ & {Non} \\ \hline
		$x + 2y = 3$ & {$\implies$} & $x^2 + 2yx = 3x$ & {$\times x$} & {Non} \\ \hline
		$x-3y  =0$ & {$\iff$} & $-2x+6y = 0$ & {$\times(-2)$ ou $\times\frac{-1}2$} & {Oui} \\ \hline
		$-2x + 3y = -3$ & {$\centernot\iff$} & $4x - 6y = -6$ & & {Non} \\ \hline
		$x^2 = x$ & {$\impliedby$} & $x = 1$ & & {Non} \\ \hline
		$x = 3$ & {$\implies$} & $x^2 = 9$ & {Mise au carré} & {Non} \\ \hline
	\end{tabular}
	\end{center}
}

\exe{difficulty=1}{
	Un élève de Seconde fournit la production écrite suivante.
	Celle-ci présente quelques erreurs.
	
	\fbox{\parbox{\textwidth}{
		\textbf{\oldunderline{Énoncé :}}
		
		Considérons l'équation $x^3 + x^2 + x + 1 = 0$ où $x\in\R$ est un nombre réel encore inconnu.
		\begin{enumerate}
			\item
			Montrer que $(x-1)(x^3 + x^2 + x + 1) = x^4 - 1$.
			\item
			En déduire la ou les solutions de l'équation originelle.
		\end{enumerate}
		
		\textbf{\underline{Réponse de l'élève :}}
		
		\begin{enumerate}
			\item
			Pour développer $(x-1)(x^3 + x^2 + x + 1)$,
			je distribue d'abord le $x$, puis le $-1$, et j'ajoute les résultats.
				\[ x(x^3 + x^2 + x + 1) + (-1)(x^3 + x^2 + x + 1) = x^4 + x^3 + x^2 + x -x^3 -x^2 - x - 1. \]
			Les puissances de $x$ égales se regroupent pour simplifier l'expression
				\[ x^4 + x^3 + x^2 + x -x^3 -x^2 - x - 1 = x^4 - 1, \]
			comme voulu.
			\item
			Un $x\in\R$ pour lequel $x^3 + x^2 + x + 1 = 0$ vérifie donc $0 = x^4 - 1$ et donc
			$x^4  = 1$.
			En appliquant la racine carrée, je trouve $x^2 = 1$ car $x^4 = (x^2)^2$.
			En l'appliquant encore, j'obtiens finalement $x = 1$.
			
			L'unique solution de l'équation est donc $1$.
		\end{enumerate}
	}}
	\begin{enumerate}
		\item 
		Identifier les erreurs commises par l'élève.
		\item
		Répondre à l'énoncé.
	\end{enumerate}
}{exe:erreurs0}{
	\begin{enumerate}
		\item
		La première question est bien traitée et ne comporte aucune erreur.	
		\item
		La deuxième question contient deux erreurs : une lors de la résolution de $x^4 = 1$, et l'autre dans sa conclusion.
		
		Premièrement, appliquer la racine carrée est pertinente, mais $\sqrt{x^2} = |x|$, et non pas $x$ !
		La valeur absolue est importante pour ne pas oublier de solution.
		Ainsi, si $x^4 = 1$, alors $|x^2| = 1$. Comme $x^2 \geq 0$ pour tout $x\in\R$, la valeur absolue peut disparaître ici.
		Puis, si $x^2 = 1$, alors $|x| = 1$, et donc $x=1$ ou $x=-1$. Il y a donc deux solutions réelles à l'équation $x^4  = 1$.
		
		Deuxièmement, la conclusion est erronnée car elle présente une erreur de logique.
		En effet, l'impliquation $x^3 + x^2 + x + 1 = 0 \implies x \in \{-1 ; 1\}$ est vraie, mais l'implication inverse ne l'est pas !
		Nous disposons de deux candidats de solutions qu'il faut vérifier. La solution $x=1$ ne fonctionne pas, mais $x=-1$, elle, fonctionne.
		La seule solution est donc $x=-1$.
		
		En multipliant l'équation originelle par $x-1$, nous avons ajouté artificiellement la solution $x=1$.
		Par exemple, les implications suivantes sont vraie
			\[ x - 1 = 0 \implies (x-3)(x-1) = 0 \implies x \in \{3 ; 1\}. \]
	\end{enumerate}
}

\exe{}{
	Résoudre les équations suivantes pour $x\in\R$ non nul. 
	%Exprimer $x$ sous la forme $q \sqrt{n}$ avec $q\in\Q$ rationnel et $n\in\N$ le plus petit possible.
	\begin{multicols}{4}
	\begin{enumerate}
		\item $\frac{x}3 = \frac12.$
		\item $ \frac7x  =\frac12$
		\item $\frac6x  =\frac{\sqrt{3}}2$
		\item $\frac{x}2 =\frac1{\sqrt{3}}$
		%\item $\frac9x  =\frac{\sqrt{6}}5$
		%\item $-\frac3x =\sqrt{7}$
		%\item $-\frac7{2x}  =\frac{\sqrt{11}}3$
		%\item $\frac3{5x} =-\sqrt{5}$
	\end{enumerate}
	\end{multicols}
}{exe:eq-inv-x}{
	\begin{multicols}{2}
	\begin{enumerate}
		\item 
			\begin{align*}
				\frac{x}3 &= \frac12 \\
				x &= 3 \cdot \frac12 \\
				x &= \frac32
			\end{align*}
		\item
			\begin{align*}
				\frac7x  &=\frac12 \\
				7 &= x \cdot \frac12 \\
				x &= 7\cdot2 = 14
			\end{align*}
		\item
			\begin{align*}
				\frac6x  &= \frac{\sqrt{3}}2 \\
				6 &= x \cdot \frac{\sqrt{3}}2 \\
				x &= 6\cdot\frac2{\sqrt{3}} \\
				x &= \frac{12}3 \sqrt3 = 4\sqrt3
			\end{align*}
		\item
			\begin{align*}
				\frac{x}2 &= \frac1{\sqrt{3}} \\
				x &= 2 \cdot \frac1{\sqrt{3}} \\
				x &= \frac23 \sqrt3
			\end{align*}
	\end{enumerate}
	\end{multicols}
}

\exe{}{
	Trouver le $x \in \Q$ rationnel (et tel que chaque dénominateur est non nul) vérifiant chaque équation suivante, et l'exprimer sous forme de fraction irréductible.
	
	\begin{multicols}{3}
	\begin{enumerate}[itemsep=10pt]
		\item $9 = \frac3x$
		\item $-4 = \frac2x$
		\item $0,2422 = \frac1x$
		\item $0,7 = \frac{0,1}x$
		\item $3 - \frac4x = -9 + \frac1x$
		\item $-2 + \frac{-2}x = \frac3x - 10$
		\item $\frac1{x-2} = 3$
		\item $\frac3{x+5} = \frac7{x+5} - 3$
	\end{enumerate}
	\end{multicols}

}{exe:eq-lin0}
{
	\begin{multicols}{2}
	\begin{enumerate}[itemsep=10pt]
		\item $9 = \frac3x \iff x = \frac13$
		\item $-4 = \frac2x \iff x=-\frac12$
		\item $0,2422 = \frac1x \iff x = \frac{10000}{2422} = \frac{5000}{1211}$
		\item $0,7 = \frac{0,1}x \iff x = \frac17$
		\item $3 - \frac4x = -9 + \frac1x \iff x = \frac5{12}$
		\item $\frac1{x-2} = 3 \iff x=\frac73$
		\item $\frac3{x+5} = \frac7{x+5} - 3 \iff x=\frac{-11}3$
	\end{enumerate}
	\end{multicols}
}


\exe{difficulty=1}{
	Montrer que les nombres suivants sont rationnels en les exprimant sous forme de fraction d'entiers.
	\[
	\begin{aligned}
		A &= 0,666{\dots} \\
		B &= 1,666{\dots} \\
	\end{aligned}
	\hspace{5cm}
	\begin{aligned}
		C &= 0,121212{\dots} \\
		D &= 0,34777{\dots} \\
	\end{aligned}
	\]
	\[
	E = 0,123456789123456789123{\dots} \text{ (nombre d'Amandine) }
	\]
	
	\emph{Indication : montrer d'abord que $10A = 6+A$.}
}{exe:dev-to-fraction}{
	\begin{enumerate}[label=\Alph*.]
		\item 
		La relation $10A = 6 + A$ donne $A = \frac23$.
		\item
		La relation $10B = 15 + B$ donne $B = \frac{15}{9} = \frac53$.
		On aurait aussi pû utiliser que $B = 1+A$.
		\item
		La relation $100C = 12 + C$ donne $C = \frac{12}{99} = \frac{4}{33}$.
		\item
		En posant $\tilde{E} = 100E = 34,777{\dots}$, on obtient $10\tilde{E} = 313 + \tilde{E}$, d'où $\tilde{E} = \frac{313}{9}$, et donc $E = \frac{1}{100}E' = \frac{313}{900}$.
		\item 
		La relation $10^9 E = 123~456~789 + E$ donne $E = \frac{123~456~789}{999~999~999}$.
	\end{enumerate}
}

\exe{}{
	Montrer que les nombres suivants sont rationnels en les exprimant sous forme de fraction d'entiers.
	\begin{align*}
		A = 0,123123123{\dots} &&
		B = 0,9999{\dots}
	\end{align*}
}{exe:dev-to-fractionE}{
	\begin{enumerate}[label=\Alph*.]
		\item
		La relation $1000A = 123 + A$ donne $A = \frac{123}{999} = \frac{41}{333}$.
		\item 
		La relation $10B = 9 + B$ donne $9B = 9$ et $B=1$.
	\end{enumerate}
}


\exe{, difficulty=1}{
	Multipliez l'équation 
		\[ \frac{a}b + \frac{c}d = x \]
	par $b$ puis $d$ à gauche et à droite et en déduire la règle d'addition des fractions :
		\[ x = \frac{ad + cb}{bd}. \]
}{exe:common-denom}{
	En multipliant par $bd$, on trouve
		\[ (bd) x =  bd \left( \frac{a}b + \frac{c}d \right) = ad + bc, \]
	d'où
		\[ x = \frac{ad+bc}{bd}. \]
}




%%%%%%%%%%%%%%%%%%%%%%%%
%%%%%%%%%% SET 3 %%%%%%%%%%
%%%%%%%%%%%%%%%%%%%%%%%% 


\exe{}{
	On souhaite connaître les solutions de l'équation du deuxième degré d'inconnue $x\in\R$ :
		\begin{align}
			22x^2 - 125x + 22  = 0. \label{eq:1}
		\end{align}
	\begin{enumerate}
		\item Montrer que $22x^2 - 125x + 22 = (2x-11)(11x-2)$.
		\item En déduire l'ensemble des solutions de l'équation \eqref{eq:1}.
	\end{enumerate}
}{exe:eq7}{
	\begin{enumerate}
		\item On part toujours de la forme factorisée qu'on développe par double distributivité :
			\[ (2x - 11)(11x - 2) = 22x^2 - 4x - 121x + 22 = 22x^2 - 125x + 22. \]
		\item On utilise que le produit est nul si et seulement si un des deux facteurs est nul.
		L'ensemble des solutions de \eqref{eq:1} est donc $\left\{ \frac{11}2 ; \frac2{11} \right\}$.
	\end{enumerate}
}


\exe{}{
	Pour chaque propriété, donner une fonction polynomiale $f$ à coefficients réels et non nulle la vérifiant.
	\begin{multicols}{2}
	\begin{enumerate}
		\item $f$ s'annule en $1$.
		\item $f$ s'annule en $-10$.
		\item $f$ s'annule en $0$.
		\item $f$ s'annule en $-10$ et en $1$.
		\item Les racines de $f$ sont $2, -3, \frac27$, et $0$.
	\end{enumerate}
	\end{multicols}
}{exe:f-zeroes}{
	Le polynôme identiquement nul $f(x) = 0$ n'est bien sûr pas autorisé sinon l'exercice n'a pas d'intérêt !
	\begin{enumerate}
		\item $f(x) = x-1$ fonctionne ainsi que $f(x) = 3(x-1)$, ou $f(x) = -(x-1) = 1-x$.
		\item $f(x) = x+10$ fonctionne, ainsi que $f(x) = (x+10)(x-1)$, et tous ses multiples.
		\item $f(x) = x$ ou $f(x) = 500x$.
		\item $f(x) = (x+10)(x-1) = x^2 +9x -10$.
		\item $f(x) = (x-2)(x+3)(x-\frac27)x$.
	\end{enumerate}

}

\exe{}{
	Résoudre les équations suivantes pour $x\in\R$ non nul. Exprimer $x$ sous la forme $q \sqrt{r}$ avec $q\in\Q$ rationnel et $r\in\N$ le plus petit possible.
	\begin{multicols}{2}
	\begin{enumerate}
		\item $\frac3x = \frac73.$
		\item $ \frac7x  =\frac1x + 2$
		\item $\frac{16}x  =\frac{\sqrt{3}}2$
		\item $-\frac2x =4\sqrt{3}$
		\item $\frac9{-x}  =\frac{\sqrt{6}}5$
		\item $-\frac3x =\frac57\sqrt{7}$
		\item $-\frac9{2x}  =\frac{\sqrt{11}}3$
		\item $\frac2{7x} =-\sqrt{14}$
	\end{enumerate}
	\end{multicols}
}{exe:eq-sqrt}{
	Nous avons, pour $x\neq0$,
		\begin{align*}
			 \frac{a}x &= \frac{b}c \\
			  \frac{a}x \times x &= \frac{b}c \times x \\
			  a &= \frac{b}c x \\
			  x &= a \times \frac{c}{b} = \frac{a\times c}{b}
		\end{align*}
	non sans rappeler le produit en croix !
	Les équations sont équivalentes ici car $x$ est supposé non nul ce qui rend la multiplication par $x$ une opération réversible.
	
	De plus, nous utilisons que 
		\[ \frac1{\sqrt{d}} = \frac{1\times\sqrt{d}}{\sqrt{d} \times \sqrt{d}} = \frac1d \sqrt{d} \]
	pour obtenir la forme requise par l'énoncé.
	\begin{multicols}{2}
	\begin{enumerate}[itemsep=10pt]
		\item
		$x = \frac{3\times3}{7} = \frac97$.
		\item
		$\frac6x = 2$, donc $ x = \frac62 = 3$.
		\item
		$ x = \frac{16\times2}{\sqrt3} = \frac{32}{\sqrt3} = \frac{32}{3} \sqrt3$.
		\item
		$ x = \frac{-2}{4\sqrt3} = \frac{-1}6 \sqrt3$.
		\item
		$x = \frac{9\times5}{\sqrt6} = \frac{15}{2}\sqrt6$.
		\item
		$ x = \frac{-3\times7}{5\sqrt7} = \frac{-3}5 \sqrt7$.
		\item
		$x = \frac{-9\times3}{2\sqrt{11}} = \frac{-27}{22}\sqrt{11}$.
		\item
		$x = \frac{2}{-7\sqrt{14}} = \frac{-1}{49}\sqrt{14}$
	\end{enumerate}
	\end{multicols}
}




\exe{}{
	On souhaite connaître les solutions de l'équation du troisième degré d'inconnue $x\in\R$ :
		\begin{align}
			4x^3 - 6x^2 - 2x + 3  = 0. \label{eq:2}
		\end{align}
	\begin{enumerate}
		\item Montrer que $4x^3 - 6x^2 - 2x + 3 = (2x-3) \left(2x^2-1\right)$.
		\item En déduire l'ensemble des solutions de l'équation \eqref{eq:2}.
	\end{enumerate}
}{exe:eq8}{

	\begin{enumerate}
		\item On part toujours de la forme factorisée qu'on développe par double distributivité.
		On utilise aussi que $x^2 \cdot x = x \cdot x \cdot x = x^3$.
			\[ (2x - 3)(2x^2 - 1) = 4x^3 - 2x -6x^2 + 3. \]
		\item On utilise que le produit est nul si et seulement si un des deux facteurs est nul.
		Or l'équation $2x^2 - 1 = 0$ est équivalente à $x^2 = \frac12$.
		En prenant la racine et en utilisant que $\sqrt{x^2} = |x|$, on obtient
			\[ |x| = \sqrt{\frac12} = \frac1{\sqrt2}. \]
		Les deux valeurs $\pm \frac1{\sqrt2}$ (notation $\pm$ pour \og plus ou moins \fg) sont donc solutions.
		On en déduit que l'ensemble des solutions de \eqref{eq:2} est $\left\{ \frac32 ; \pm \frac1{\sqrt2} \right\} = \left\{ \frac32 ; \frac1{\sqrt2} ; - \frac1{\sqrt2} \right\}$
	\end{enumerate}

}


\exe{}{
	On souhaite connaître les solutions de l'équation du quatrième degré d'inconnue $x\in\R$ :
		\begin{align}
			16x^4 - 24x^2 + 9  = 0. \label{eq:3}
		\end{align}
	\begin{enumerate}
		\item Montrer que $16x^4 - 24x^2 + 9 = \left(4x^2-3\right)^2$.
		\item En déduire l'ensemble des solutions de l'équation \eqref{eq:3}.
	\end{enumerate}
}{exe:eq9}{

	\begin{enumerate}
		\item On part toujours de la forme factorisée qu'on développe à l'aide de l'identité remarquable $(a-b)^2 = a^2 + b^2 -2ab$.
		On utilise aussi que $(x^2)^2 = x^2 \cdot x^2 = x\cdot x\cdot x\cdot x = x^4$.
			\[ (4x^2 - 3)^2 = (4x^2)^2 + 3^2 - 2\cdot(4x^2)\cdot3 = 16x^4 + 9 -24x^2. \]
		\item On utilise que le produit est nul si et seulement si un des deux facteurs est nul.
		Ici, les deux facteurs sont les mêmes, car $(4x^2 - 3)^2 = (4x^2 - 3)(4x^2 - 3)$, donc on se remet à résoudre $(4x^2 - 3) = 0$.
		Ceci est équivalent à $x^2 = \frac34$, et donc l'ensemble des $x$ vérifiant \eqref{eq:3} est donné par $\left\{ \pm \sqrt{\frac34} \right\} = \left\{ \pm \frac12\sqrt{3} \right\}= \left\{ \frac12\sqrt{3} ;  -\frac12\sqrt{3} \right\}$.
	\end{enumerate}

}


\exe{difficulty=1}{
	\begin{multicols}{2}
	La Terre complète son tour autour du Soleil en environ 365,2422 jours. 
	Un calendrier ne peut donc pas contenir 365 jours tous les ans si celui-ci doit être synchronisé avec les saisons.
	
	Si le calendrier contient moins de 365,2422 jours, il sera en avance.
	S'il contient plus de 365,2422 jours, il sera en retard

	\begin{center}
		\includegraphics[page=4]{../../CM/2-figures.pdf}
	\end{center}
	\end{multicols}
	
	\begin{enumerate}
	\item
	Calendrier julien\footnote{De Jules César (né en 100 av. J.-C., mort en 44 av. J.-C.).}.
	
	En $46$ av. J.-C., Jules César introduit la règle de l'année bissextile.
	\begin{enumerate}
		\item
		Trouver un entier naturel $a \in \N$ non nul tel que
			\[ 365,2422 \approx 365 + \frac1a.\]

		\item
		En déduire la règle d'année bissextile du calendrier julien.
		
		\item
		Au bout de 1 600 ans, de combien de jours le calendrier est-il en retard ?
	\end{enumerate}
	
	\item
	Calendrier grégorien\footnote{Du pape Grégoire XIII (1502-1585).}.
	
	En 1582, plus de 1 600 ans après l'introduction de la règle d'année bissextile, le calendrier julien prend donc un retard notable.
	Le pape Grégoire XIII promulgue alors une réforme du calendrier, ainsi créant le calendrier grégorien encore utilisé aujourd'hui.
	
	Afin de corriger le problème de dérive du calendrier, une nouvelle règle est ajoutée à la définition d'année bissextile.
	
	\begin{enumerate}
		\item
		Trouver un entier naturel $b \in \N$ tel que
			\[ 365,2422 \approx 365 + \frac14 - \frac{b}{400}.\]

		\item
		En déduire la règle d'année bissextile du calendrier grégorien.
		
		\item
		Au bout de combien d'années le calendrier est-il en retard d'un jour ?
	\end{enumerate}
	
	\item
	Un nouveau calendrier ?
	\begin{enumerate}
		\item
		Trouver un entier naturel $c \in \N$ tel que
			\[ 365,2422 \approx 365 + \frac14 - \frac{3}{400} - \frac{c}{4000}.\]

		\item 
		Créer une nouvelle règle d'année bissextile.
		\item
		Calculer le nombre d'années avant que votre calendrier soit en retard d'un jour.
	\end{enumerate}
	\end{enumerate}
}{exe:calendrier}{
	\begin{enumerate}
		\item
		\begin{enumerate}
			\item
				\[ 365,2422 \approx 365 + \frac1a \iff 0,2422 \approx \frac1a \iff a \approx \frac1{0,2422} \approx 4,13. \]
			Ainsi $a=4$ est l'entier naturel le plus proche de résoudre l'équation.
			\item
			Pour qu'une année ait $365 + \frac14$ jours en moyenne, il faut qu'elle ait $365$ jours et $366$ jours une fois tous les quatre ans.
			C'est l'année bissextile : les ans multiples de 4 ont un jour en plus, le 29 février.
			C'est le cas de l'année 2028, par exemple.
			
			\item
			L'erreur commise est égale à $365,2422 - (365 + \frac14) = -0.0078$.
			Il y a donc trop de jours dans le calendrier qu'il n'en faudrait, exactement $0,0078$ jours en trop.
			Après 1 600 ans, il y aura donc un décalage de $1600 \times 0,0078 = 12,48$ jours.
			Un tel décalage se remarque dans l'arrivée des saisons servant à l'organisation des temps de récoltes.
		\end{enumerate}
		\item
		\begin{enumerate}
			\item
				\[ 365,2422 \approx 365 + \frac14 - \frac{b}{400} \iff 0,0078 \approx \frac{b}{400} \iff b \approx 400\times0,0078 = 3,12. \]
			Ainsi $b=3$ est l'entier naturel le plus proche de résoudre l'équation.
			\item
			Pour qu'une année ait $365 + \frac14 - \frac3{400}$ jours en moyenne, il faut garder la règle de l'année bissextile mais soustraire trois jours tous les 400 ans.
			Une façon de faire est de soustraire un jour du calendrier tous les 100 ans, sauf les multiples de 400.
			Une année est donc dite bissextile si elle est multiple de 4 sauf multiple de 100 mais pas de 400.
			
			\item
			L'erreur commise est égale à $365,2422 - (365 + \frac14 - \frac3{400}) = -0.0003$.
			Il y a donc trop de jours dans le calendrier qu'il n'en faudrait, exactement $0,0003$ jours en trop.
			Au bout de $n$ années, il y aura $0,0003 \times n$ jours en trop.
			Pour atteindre 1 jour en trop il faut donc que $0,0003 \times n \approx 1 \iff n \approx \frac1{0,0003} \approx 3~333,33$.
			
			En conclusion, après 3 333 années du calendrier grégorien, celui-ci aura accumulé 1 jour de retard.
		\end{enumerate}
		\item
		\begin{enumerate}
			\item
			On trouve $c=1$.
			\item
			Il faut soustraire un jour du calendrier tous les 4 000 ans.
			Une année bissextile pourrait donc être tous les multiples 4, sauf les multiples de 100 pas multiples de 400, et sauf multiples de 4 000.
			\item 
			Après un an, l'erreur commis est de $5\times10^{-5} = 0,00005$ jours en trop dans le calendrier.
			Il faut donc $\frac1{5\times^{-5}} = 0,2 \times 10^5 = 20~000$ années pour que ce calendrier soit en retard d'un jour.
		\end{enumerate}
	\end{enumerate}
}





\exe{, difficulty=2}{
	Donner quelques termes en plus de la somme infinie suivante. 
		\[ \frac1{10} + \frac1{100} + \frac1{1 000} + \frac1{10 000} + \dots \]
	Écrire le nombre en écriture décimale, puis l'exprimer sous forme de fraction de deux entiers comme à l'exercice \ref{exe:dev-to-fraction}.
}{exe:20}{
	La somme infinie vaut $0,1111{\dots} = \frac19$.
}

\exe{, difficulty=2}{
	Donner quelques termes en plus de la somme infinie suivante.
		\[ \frac12 + \frac14 + \frac18 + \frac1{16} + \dots \]
	À quoi est-elle égale ? Un résultat entier est attendu.
}{exe:21}{
	Posons $S$ cette somme. 
	En multipliant par $2$, on trouve $2S = 1 + S$, et donc $S=1$.
}

\exe{, difficulty=2}{
	On considère la somme finie suivante
		\[ S = 1 + 2 + 4 + 8 + 16 + \dots + 2^{n-1} +  2^n, \]
	où $n\in\N$ est un entier naturel.
	
	Montrer que $2S = S - 1 + 2^{n+1}$, et en déduire la valeur de $S$ sous forme close (sans points de suspension).
}{exe:22}{
	En multipliant par 2, le premier terme (1) disparait, et un nouveau terme ($2^{n+1}$) apparaît.
	On a donc bien $2S = S -1 + 2^{n+1}$, et donc $S = 2^{n-1} - 1$.
}


\exe{, difficulty=1}{
	Montrer qu'un triangle isocèle rectangle ne peut pas avoir des côtés de longueurs toutes entières.
}{exe:non-iso-rec}{
	Supposons qu'il existe un triangle isocèle rectangle de côtés $(a ; a; b)$.
	
	Or, comme le triangle est rectangle isocèle,
		\[ b^2 = a^2 + a^2 = 2a^2 \iff b = \pm\sqrt2 a. \]
	Par conséquent, $\sqrt2 = \pm\frac{b}{a} \in \Q$, ce qui contredit le théroème d'irrationalité de la racine carrée de 2.
}

\exe{, difficulty=1}{
	Soit $q\in\Q$ un rationnel. Montrer qu'un triangle dont les côtés sont de longueur $(4 ; 5 ; q)$ ne peut être rectangle que si $q=3$.
	
	En autorisant $x\in\R$ irrationel, donner un triangle rectangle de côtés $(4 ; 5; x)$ avec $x\neq3$.
}{exe:non-rec}{
	Il y a deux cas de figure à étudier : $q$ est la longueur de l'hypothénuse, ou elle ne l'est pas.
	Dans le premier cas, le théorème de Pythagore énonce que
		\[ 4^2 + 5^2 = q^2 \iff 41 = q^2. \]
	Il suit que $q=\sqrt{41}$ qui n'est pas rationnel car 41 n'est pas un carré parfait.
	Ceci répond à la deuxième question : $(4 ; 5; \sqrt{41})$ est rectangle.
	
	Dans le second cas,
		\[ q^2 + 4^2 = 5^2 \iff q^2 = 9. \]
	Par conséquent, $q = \pm3$ et seule la solution $q=3$ est positive et donc une longueur valide.
}



%%%%%%%%%%%%%%%%%%%%%%%%
%%%%%%%%%% SET 4 %%%%%%%%%%
%%%%%%%%%%%%%%%%%%%%%%%% 




\exemulticols{}{
	Donner approximativement les coordonnées de chaque point du repère ci-contre.
	\begin{align*}
		&A( \qquad ; \qquad) \\[10pt]
		&B \\[10pt]
		&C \\[10pt]
		&D \\[10pt]
		&E
	\end{align*}
	\vfill\null
}{
	\begin{center}
	\includegraphics[page=5]{../../CM/2-figures.pdf}
	\end{center}
}{exe:lecture-coord}{
	\begin{align*}
		A(3 ; 0) && B(-1 ; 3) && C(3 ; 3) &&
		D(0;-2) && E(-2,5 ; 0)
	\end{align*}
}



\exe{, difficulty=0}{
	En reprenant les points de l'exercice \ref{exe:lecture-coord}, calculer
		\begin{enumerate}
			\item $BC$, la longueur du segment $[BC]$ ;
			\item $AC$, la longueur du segment $[AC]$ ; et
			\item $AB$, la longueur du segment $[AB]$.
		\end{enumerate}
}{exe:longueur-ex}{
	\begin{enumerate}
		\item Les points sont alignés horizontalement, et on lit que $BC=4$, différence des abscisses.
		\item Les points sont alignés horizontalement, et on lit que $AC=3$, différence des ordonnées.
		\item Le triangle $ABC$ est rectangle. On a donc $AB^2 = AC^2 + AC^2 = 25$ d'après le théorème de Pythagore. 
		Il suit que $AB =5$.
	\end{enumerate}
}

\exe{difficulty=1}{
	Généraliser l'exercice \ref{exe:longueur-ex} en donnant la formule de la longueur d'un segment dans un repère orthonormé.
	Pour $A(x_A ; y_A)$ et $B(x_B ; y_B)$ deux points,
		\[ AB^2 = \phantom{(x_B-x_A)^2 + (y_B-y_A)^2.}\]
}{exe:formule-longueur}{
	Le théorème de Pythagore donne 
		\[ AB^2 = (x_B-x_A)^2 + (y_B-y_A)^2.\]
}



\exe{}{
	On considère le triangle de sommets $A(1;2), B(-3 ; 5), C(-5 ; -6)$.
	\begin{enumerate}
		\item Tracer le triangle $ABC$ dans un repère.
		\item Calculer le carré de chacun des côtés. 
		\item Que dire du triangle ?
	\end{enumerate}
}{exe:Trex}{
	\begin{multicols}{2}
	\begin{center}
	\includegraphics[page=6]{../../CM/2-figures.pdf}
	\end{center}
	
	\begin{enumerate}[start=1]
		\item
			\begin{align*}
				AB^2 &= \norm{A-B}^2 = \norm{(4 ; -3)}^2 = 16 + 9 = 25, \\
				AC^2 &= \norm{A-C}^2 = \norm{(6 ; 8)}^2 = 36 + 64 = 100, \\
				BC^2 &= \norm{B-C}^2 = \norm{(2 ; 11)}^2 = 4 + 121 = 125.
			\end{align*}
		\item
		D'après la réciproque du théorème de Pythagore, le triangle est rectangle en $A$ car $BC^2 = AB^2 + AC^2$.
	\end{enumerate}
	\end{multicols}
}

\exe{}{
	Considérons les points $\point{A}{1}{1}, \point{B}{3}{1}, C\bigl(2 ; \sqrt{3}+1\bigr)$.
	Démontrer que le triangle $ABC$ est équilatéral en calculant le carré de la longueur de chaque côté.
	
	%\emph{Un triangle équilatéral est un triangle dont les trois côtés ont la même longueur}.
}{exe:equilateral}{
	On calcule 
		\begin{align*}
			AB^2 = {2^2 + 0^2} = 4, && AC^2 = {1^2 + \sqrt{3}^2} = 4, && BC^2 = 4.
		\end{align*}
}

\exe{, difficulty=1}{
	Soient $A(1;1), B(3;1)$ deux points du plan, et $C(2;x)$ un point dépendant d'un paramètre réel $x\in\R$.
	
	\begin{enumerate}
		\item Où se situe le point $C$ ? Tracer dans un repère les points $A, B$, ainsi que les positions possibles de $C$.
		\item
		Pour quel(s) $x\in\R$ le triangle $ABC$ est-il équilatéral ? 
	\end{enumerate}
}{exe:equilateral2}{
	On souhaite que $AB = AC = BC$ soit vérifié, c'est-à-dire que
		\[ 2 = \sqrt{1 + (1-x)^2} = \sqrt{1 + (1-x)^2}. \]
	La deuxième égalité est redondante et, en mettant au carré, on trouve
		\begin{align*}
		4 = 1 + (1-x)^2 && \iff && (1-x)^2 = 3.
		\end{align*}
	En reprenant l'exercice \ref{exe:carre4-all}, on trouve deux solutions :
	\vspace{-20pt}
		\begin{multicols}{2}
		\begin{align*}
			1-x &= \sqrt{3}, \\
			x &= 1-\sqrt{3}.
		\end{align*}
			
		\begin{align*}
			1-x &= -\sqrt{3}, \\
			x &= 1 +\sqrt{3}.
		\end{align*}
		\end{multicols}
}{}

\exe{}{
	Considérons les points
		\begin{align*}
			\point{A}{0}{1}, && \point{B}{-3}{0}, && \point{C}{1}{-2}, && \point{D}{-2}{-3}.
		\end{align*}
	Démontrer que le quadrilatère $BACD$ est un parallélogramme en comparant le milieu de ses deux diagonales.
	
	\emph{Un parallélogramme est un quadrilatère dont les diagonales se coupent en leur milieu}.
}{exe:parallelogramme}{
	Le milieu du segment $[BC]$ est donné par
		\[ M = \frac12 (B+C) = (-1;-1),\]
	et le milieu du segment $[AD]$ est donné par
		\[ M' = \frac12 (A+D) = (-1;-1) = M.\]
	Le quadrilatère est donc bien un parallélogramme.

	\centering
	\includegraphics[page=7]{../../CM/2-figures.pdf}
}

\exe{, difficulty=1}{
	Soient $\point{A}12$ et $\point{M}3{-1}$ deux points du plan.
	Quel point $C$ faut-il choisir pour que $M$ soit le milieu du segment $[AC]$ ? Donner ses coordonnées.
}{exe:milieu2}{
	La contrainte que $M$ soit le milieu de $[AC]$ s'écrit
		\begin{align*}
			M = \frac12(A+C) && \iff && 2M = A+C && \iff && C = 2M-A
		\end{align*}
	Par suite,
		\[ C= 2M-A = 2\cdot\left(3;-1\right) - \left(1;2\right) = \left(5; -4\right). \]
	
}


\exe{, difficulty=1}{
	Considérons $A(-2;-1), B(-1;1)$, et $C(-3 ; 2)$ trois points.
	Quel point $D$ poser de telle sorte que le quadrilatère $ABDC$ soit un parallélogramme ?
	
	\emph{Indication : calculer le milieu $M$ de $BC$. $M$ doit aussi être le milieu de $AD$.}
}{exe:parall-constr}{

	\begin{multicols}{2}
	\begin{center}
	\includegraphics[page=8]{../../CM/2-figures.pdf}
	\end{center}
	
	Calculons d'abord 
		\[ M = \frac{B+C}2 = \left(-2 ; \frac32 \right). \]
	
	Imposons ensuite $M  = \frac{A+D}2$, équivalent à $D = 2M - A$.
	Il suit que 
		\[ D = (-4 +2 ; 3 +1) = (-2 ; 4). \]
	\end{multicols}
}


\exe{, difficulty=1}{
	Soient $p, q \in \Z$ deux entiers.
	Montrer que le point $\bigl(p^2 - q^2 ; 2pq\bigr)$ est à distance $p^2 + q^2$ de l'origine.
}{exe:norme-triplet-pythagoricien}{
	On utilise que $(a+b)^2 = a^2 + b^2 + 2ab$ et que $(a-b)^2 = a^2 + b^2 - 2ab$ pour calculer la norme du point.
	Ces deux identités, dites \emph{remarquables}, seront étudiées plus tard dans le cours.
	Nonobstant, la double distributivité permet de les démontrer facilement.
		\begin{align*}
			\bigl(p^2 - q^2\bigr)^2 + (2pq)^2 &= p^4 + q^4 - 2p^2q^2 + 4p^2q^2, \\
												&= p^4 + q^4 + 2p^2q^2, \\
												&= \bigl(p^2 + q^2\bigr)^2.
		\end{align*}
}


%%%%%%%%%%%%%%%%%%%%%%%%
%%%%%%%%%% SET 5 %%%%%%%%%%
%%%%%%%%%%%%%%%%%%%%%%%% 


\exe{, difficulty=0}{
	Un étudiant jette une balle dans les airs et mesure la hauteur de la balle tous les quarts de seconde.
	Il note ses résultats dans le tableau ci-dessous.
	\begin{center}
		\def\arraystretch{1.5}
		%\setlength\tabcolsep{20pt}
		\begin{tabular}{|c|c|c|c|c|c|c|c|}\hline
			Hauteur (cm) & 85 & 145 & 190 & 145 & 85 & 40 & 0 \\ \hline
			Temps (s) & 0 & 0,25 & 0,5 & 0,75 & 1 & 1,25 & 1,5 \\\hline
		\end{tabular}
	\end{center}
	
	\begin{enumerate}
		\item La hauteur est-elle une fonction du temps ? Justifier.
		\item Le temps est-il une fonction de la hauteur ? Justifier.
	\end{enumerate}
}{exe:f1}{
	\begin{enumerate}
		\item On peut associer \textbf{une seule} hauteur à chaque temps. La hauteur peut donc être vue comme fonction du temps.
		\item On ne peut pas associer un unique temps à chaque hauteur. Par exemple, il y a deux temps distincts pour lesquels la hauteur est de $85$cm (0s et 1s).
		Le temps ne peut donc pas être vu comme fonction de la hauteur.
	\end{enumerate}
}

\exe{, difficulty=0}{
	Montrer que
	\begin{enumerate}
		\item
		le rayon $R$ d'un cercle est fonction de son périmètre $P$ et écrire la fonction $\text{Rayon}(P)$ associée ;
		\item
		l'aire $A$ d'un cercle est fonction de son rayon $R$ et écrire la fonction $\text{Aire}(R)$ associée ; et
		\item
		l'aire $A$ d'un cercle est fonction de son périmètre $P$ et écrire la fonction $\text{Aire}(P)$ associée.
	\end{enumerate}
}{exe:f2}{
	\begin{enumerate}
		\item
		De la relation 
			\[ P = 2\pi \cdot R, \]
		on déduit que 
			\[ R = \frac{1}{2\pi} \cdot P. \]
		Remarquons qu'on a inversé le membre de gauche et le membre de droite pour extraire la forme d'une fonction : une valeur de $P$ donne une unique valeur de $R$.
		
		Ainsi on peut voir $R$ comme une fonction de $P$. On note alors $R=R(P)$, qui vérifie
		\[ R(P) = \frac{1}{2\pi} \cdot P. \]
		
		\item 
		La formule de l'aire
			\[ A = \pi \cdot R^2 \]
		exprime directement l'aire comme fonction du rayon : à chaque rayon possible, on peut calculer une unique aire.
		On note alors
			\[ A(R) =  \pi \cdot R^2. \]
		\item
		On souhaite exprimer l'aire $A$ comme fonction du périmètre $P$.
		Or d'après les questions précédentes, l'aire $A$ est fonction du rayon $R$, et le rayon $R$ est lui-même fonction du périmètre $P$.
		
		Ainsi, une valeur de $P$ donne une unique valeur de $R$ qui donne une unique valeur de $A$.
		En suivant le raisonnement, on peut composer les fonctions trouvées ci-dessus en considérant $A(R(P))$.
		Pour ne pas surcharger les notations, appelons plutôt $\text{Aire}$ la fonction prenant un périmètre.
			\begin{align*}
			 \text{Aire}(P) &= A\bigl(R(P)\bigr) \\ &=A\left(\frac{1}{2\pi} \cdot P\right) \\ &= \pi \left( \frac{1}{2\pi} \cdot P \right)^2 \\ &= \frac{\pi}{4\pi^2} \cdot P^2 = \frac{1}{4\pi} \cdot P^2.
			 \end{align*}
	\end{enumerate}
}


\exe{}{
	Un fonction $f$ admet le tableau de valeurs suivant.
		\begin{center}
		\def\arraystretch{1.2}
		\setlength\tabcolsep{20pt}
		\begin{tabular}{|c|c|c|c|c|}\hline
			$x$ & 0 & -2 & 1 & -1 \\ \hline
			$f(x)$ & 1 & 0 & 0 & 1 \\ \hline
		\end{tabular}
		\end{center}
	Parmi les expressions algébriques suivantes, lesquelles \underline{ne peuvent pas} correspondre à $f(x)$ ?
		\begin{multicols}{4}
		\begin{enumerate}[label=\roman*)]
			\item $1-x$
			\item $1+\frac{x}2$
			\item $\frac{1-x}2$
			\item $\frac{-x^2 - x + 2}2$
		\end{enumerate}
		\end{multicols}
}{exe:f4}{
	Il s'agit de discriminer les fonctions possibles en utilisant les images du tableau.
		\begin{enumerate}[label=\roman*)]
			\item $1-x$ vaut bien $1$ en $x=0$, mais l'expression vaut $3$ en $x=-2$, donc ce n'est pas l'expression de $f$.
			\item  $1+\frac{x}2$ vaut bien $1$ en $x=0$ et $0$ en $x=-2$, mais elle vaut $1$ en $x=1$, ce n'est donc pas l'expression de $f$.
			\item $\frac{1-x}2$ vaut $\frac12$ en $x=0$, ce n'est donc pas l'expression de $f$.
			\item $f(x) = \frac{-x^2 - x + 2}2$ est la seule expression possible, qu'on vérifiera en calculant les images de $0, -2, 1,$ et $-1$.
			
			Attention aux parenthèses lors de la substitution d'une valeur négative ainsi qu'à l'ordre des opérations.
			Par exemple, en $x=-2$,
				\[ \frac{-x^2 - x + 2}2 = \frac{-(-2)^2 - (-2) + 2}2 = \frac{-(4)+2+2}2 = 0. \]
		\end{enumerate}

	Cet exercice deviendra vite évident après avoir terminé l'étude des fonctions affines. 
	Comme les courbes représentatives des trois premières expressions sont des droites non constantes, toutes leurs images ont un unique antécédent.
	Or ici $0$ et $1$ ont deux antécédents.
}



\exe{}{
	Remplir le tableau d'images suivant.
		\begin{center}
		\def\arraystretch{1.2}
		\setlength\tabcolsep{20pt}
		\begin{tabular}{|c|c|c|c|c|}\hline
			$x$ & 0 & -2 & 1 & -1 \\ \hline
			$f(x)$ & 1 & 3 & 0 & \\ \hline
			$g(x)$ & 1 & 0 &  & -2 \\ \hline
			$f(x)+2g(x)$ &  & & 2 & 1 \\ \hline
		\end{tabular}
		\end{center}
}{exe:fonctions-QCM-trous}{
	\begin{center}
	\def\arraystretch{1.2}
	\setlength\tabcolsep{20pt}
	\begin{tabular}{|c|c|c|c|c|}\hline
		$x$ & 0 & -2 & 1 & -1 \\ \hline
		$f(x)$ & 1 & 3 & 0 & 5 \\ \hline
		$g(x)$ & 1 & 0 & 1 & -2 \\ \hline
		$f(x)+2g(x)$ & 3 & 3 & 2 & 1 \\ \hline
	\end{tabular}
	\end{center}
}


\exe{, difficulty=0}{
	Posons $h(x) = (x+2)(x-1)$ pour tout $x\in\R$.
	Remplir le tableau de valeurs ci-dessous.
		\begin{center}
		\def\arraystretch{1.2}
		\setlength\tabcolsep{30pt}
		\begin{tabular}{|c|c|c|}\hline
			$x$ & -2 & 1 \\ \hline
			$f(x)$ & -42 & 1729 \\ \hline
			$h(x)$ &  &   \\ \hline
			$f(x)+h(x)$ &  &   \\ \hline
			$f(x)+2h(x)$ &  &   \\ \hline
		\end{tabular}
		\end{center}
}{exe:fonctions-QCM2}{
	Lors du calcul de chacune des images, au moins un facteur est nul : le produit vaut donc toujours zéro.
		\begin{center}
		\def\arraystretch{1.2}
		\setlength\tabcolsep{30pt}
		\begin{tabular}{|c|c|c|}\hline
			$x$ & -2 & 1 \\ \hline
			$f(x)$ & -42 & 1729 \\ \hline
			$h(x)$ & 0 & 0   \\ \hline
			$f(x)+h(x)$ & -42 & 1729  \\ \hline
			$f(x)+2h(x)$ & -42 & 1729  \\ \hline
		\end{tabular}
		\end{center}
}

\exe{, difficulty=1}{
	Montrer qu'il existe une infinité de fonctions, toutes différentes, vérifiant le même tableau de valeurs de l'exercice \ref{exe:fonctions-QCM2}.
}{exe:fonctions-QCM3}{
	On prend par exemple $f(x) + h(x), f(x) + 2h(x), f(x) + 3h(x), \dots$.
}

\exe{}{
	Considérons la fonction $f$ donnée {algébriquement} par
	\begin{align*}
		f(x) = 3x + 1
	\end{align*}
	pour tout $x\in\R$.
	
	\begin{enumerate}
		\item
		Calculer l'image par $f$ de $0$ ; de $3,1$ ; de $\frac13$ ; de $-1$ ; de $-\frac23$.
		\item
		Donner un antécédent de $1$ par $f$.
		\item
		Déterminer tous les antécédents de $6$ par $f$.	
	\end{enumerate}
}{exe:images-antecedents}{
	\begin{enumerate}
	\item
	L'image de $x$ est donnée par $f(x)$. On calcule donc
		\begin{align*}
			f(0) &= 3\cdot0 + 1 = 1 \\
			f(3,1) &= 3\cdot3,1 + 1 = 10,3 \\
			f\left(\frac13\right) &= 3\cdot\frac13 + 1 = 1 + 1 = 2 \\
			f(-1) &= 3\cdot(-1) + 1 = -2 \\
			f\left(-\frac23\right) &= 3 \cdot \left(- \frac23\right) + 1 = -1.
		\end{align*}
	\item
	La relation
		\[ f(0) = 1 \]
	implique que l'image de $0$ est $1$, et qu'un antécédent de $1$ est $0$.
	On a donc trouvé $0$, antécédent de $1$ par $f$.
	\item
	Appelons $x$ un antécédent de $6$ par $f$.
	Autrement dit, l'image de $x$ est $6$, et donc 
		\[ f(x) = 6. \]
	En utilisant l'expression algébrique de $f$, on trouve
		\begin{align*}
			f(x) &= 6, \\
			3\cdot x + 1 &= 6, \\
			3\cdot x &= 5, \\
			x &= \frac53.
		\end{align*}
	Ainsi $x=\frac53$ est le seul antécédent de $6$ par $f$.
	En général, certaines fonctions admettent plusieurs antécédents.
	Voir par exemple l'exercice suivant.
	\end{enumerate}


}






\exe{}{
	Écrire chaque intervalle suivant en termes d'inégalités qu'un réel $x \in \R$ vérifie dès qu'il appartient à celui-ci.
	\begin{multicols}{2}
	\begin{enumerate}
		\item $[-1 ; 1]$
		\item $[-3 ; 1]$
		\item $[-1 ; 1] \cap [-3 ; 1]$
		\item $]3 ; +\infty [$
		\item $] {-}\infty ; 2] \cap [-3 ; -2]$
		\item $]{-}\infty ; 4 [ \cup [2 ; +\infty[$
	\end{enumerate}
	\end{multicols}
}{exe:intervalle1}{
	\begin{multicols}{2}
	\begin{enumerate}
		\item $-1 \leq x \leq 1$
		\item $-3 \leq x \leq 1$
		\item $-1 \leq x \leq 1$
		\item $x > 3$
		\item $-3 \leq x \leq -2$
		\item $x \in \R$
	\end{enumerate}
	\end{multicols}
}

\exe{}{
	Vrai ou faux ? Justifier.
	
	\begin{multicols}{2}
	\begin{enumerate}
		\item $[-1; 3 [ \subset [-2 ; 4]$
		\item $] -6,2 ; 3] \subset [-6,1 ; 4[$
		\item $\{ -1; 3,4 \} \subset ]{-}\infty; -1] \cup [2; 5[$
		\item $ 3 \in [-3 ; 3[ \cup ]3, 6]$
	\end{enumerate}
	\end{multicols}
}{exe:intervalle2}{
	\begin{multicols}{2}
	\begin{enumerate}
		\item Vrai
		\item Faux
		\item Vrai
		\item Faux
	\end{enumerate}
	\end{multicols}

}

\exe{}{
	À quels intervalles le réel $x \in \R$ appartient-il s'il vérifie les inégalités suivantes ?
	\begin{multicols}{2}
	\begin{enumerate}
		\item $x \geq 0$
		\item $x > 1$
		\item $x \leq -3$
		\item $-3 < x \leq 8$
		\item $2x + 1 \leq 4$
		\item $-x > 0$
		\item $-3x \leq 1$
		\item $-4 \leq 2x - 1 < 9$
	\end{enumerate}
	\end{multicols}
}{exe:intervalle3}{

	\begin{multicols}{2}
	\begin{enumerate}
		\item $[0; {+}\infty[$
		\item $]1; {+}\infty[$
		\item $]{-}\infty ; -3]$
		\item $]{-3} ; 8]$
		\item $\left]{-}\infty ; \frac32\right]$
		\item $]{-}\infty ; 0[$
		\item $\left[-\frac13 ; {+}\infty\right[$
		\item $\left[-\frac32 ; 5\right[$
	\end{enumerate}
	\end{multicols}
}

\exe{}{
	À quels intervalles le réel $x \in \R$ appartient-il s'il vérifie les inégalités suivantes ?
	\begin{multicols}{2}
	\begin{enumerate}
		\item $| x - 2 | \leq 1$
		\item $| 2 + x| > 4$
		\item $| 3 + 2x | \geq 1$
		\item $| 4 - 3x | < 5$
	\end{enumerate}
	\end{multicols}
}{exe:intervalle4}{

	\begin{multicols}{2}
	\begin{enumerate}
		\item $[1 ; 3]$
		\item $]{-}\infty ; -6[ \cup ]2 ; {+}\infty[$
		\item $]{-}\infty; -2] \cup [-1; {+}\infty[$
		\item $\left] -\frac13; 3\right[$
	\end{enumerate}
	\end{multicols}
}

\exe{}{
	Écrire les unions et intersections suivantes sous forme d'intervalle.
	
	
	\begin{multicols}{2}
	\begin{enumerate}
		\item $[-1,1] \cap [-3, 1]$
		\item $]{-}\infty, 2] \cap [-3, 1]$
		\item $]{-}\infty, 4 [ \cup [2, +\infty[$
		\item $] {-}3 ; 4 [ \cap [-5 ; 4]$
		\item $]{-}2 ; 4 [  \cup [4 ; 8[$
		\item $ ]{-}\infty ; 1 [ \cap ]{-}1 ; +\infty [$
	\end{enumerate}
	\end{multicols}

}{exe:intervalle5}{
	\begin{multicols}{2}
	\begin{enumerate}
		\item $[-1,1]$
		\item $[-3, 1]$
		\item $\R$
		\item $]{-}3 ; 4[$
		\item $]{-2} ; 8[$
		\item $]{-1};1[$
	\end{enumerate}
	\end{multicols}
}

\exe{}{
	Pour chaque couples d'intervalles $I, J$, exprimer $I \cap J$ et $I \cup J$ sous forme d'intervalle.
	%\begin{multicols}{2}
	\begin{enumerate}\itemsep10pt
		\item $I = [-4 ; 2]$, et $J = [-1 ; 6[$
		\item $ I = ] {-}\infty ; 3 [$, et $J =[-2 ; {+}\infty[$ 
		\item $I = ]{-}\infty ; 7]$, et $J =]{-}\infty ; 4]$
		\item $I = ]{-}\infty ; -1]$, et $J = [-1 ; {+}\infty [$
		\item $I = [3 ; {+}\infty [$, et $J = [-5 ; {+}\infty [$
		\item $I = \{ x \in \R \ \text{tq. } x > 3 \}$, et $J = \{ x \in \R \ \text{tq. } x \leq 9 \}$
		\item $I = \{ x \in \R \ \text{tq. } 2 \leq x \leq 12 \}$, et $J = \{ x \in \R \ \text{tq. } 10 > x \geq -1 \}$
	\end{enumerate}
	%\end{multicols}
}{exe:intervalle6}{

	\begin{enumerate}\itemsep10pt
		\item $I\cap J = [{-}1; 2]$ et $I \cup J = [{-}4; 6[$
		\item $I\cap J = [{-}2; 3[$ et $I \cup J = \R$
		\item $I\cap J = ]{-}\infty; 4]$ et $I \cup J =  ]{-}\infty; 7]$
		\item $I\cap J = \{-1\} $ et $I \cup J = \R$
		\item $I\cap J [3; {+}\infty[= $ et $I \cup J = [{-}5; {+}\infty[$
		\item $I\cap J = ]3;9]$ et $I \cup J = \R$
		\item $I\cap J = [2 ; 10[$ et $I \cup J = [{-}1; 12]$
	\end{enumerate}
}


\exe{}{
	Soit $f$ définie algébriquement sur $\R^* = ]{-\infty}; 0[ \cup ]0 ; {+\infty}[$ par
		\[ f(x) = 3x^2 - 2x + 2 - \frac4x. \]
	Justifier que $f$ n'est pas définie en $0$, et calculer les images suivantes.
	
	\begin{multicols}{4}
	\begin{enumerate}
		\item $f(1)$
		\item $f(4)$
		\item $f(-1)$
		\item $f(-4)$
		\item $f(-8)$
		\item $f(8)$
		\item $f(3)$
		\item $f(-2)$
	\end{enumerate}
	\end{multicols}

}{exe:images}{
	La fonction $f$ n'est pas définie en 0 à cause de l'expression $\frac4x$ qu'il n'est pas possible d'évaluer en $x=0$, la division par 0 étant illégale.
	
	
	\begin{multicols}{4}
	\begin{enumerate}
		\item $f(1) = -1$
		\item $f(4) = 41$
		\item $f(-1) = 11$
		\item $f(-4) = 59$
		\item $f(-8) = \frac{421}2$
		\item $f(8) = \frac{355}2$
		\item $f(3) = \frac{65}3$
		\item $f(-2) = 20$
	\end{enumerate}
	\end{multicols}

}


\exe{, difficulty=1}{
	Pour chaque propriété suivante, donner algébriquement une fonction la vérifiant.
		%\begin{multicols}{1}
		\begin{enumerate}
			\item Une image admet exactement un antécédent.
			\item Une image admet exactement deux antécédents.
			\item Une image admet exactement trois antécédents.
			\item Une image admet une infinité d'antécédents.
		\end{enumerate}	
		%\end{multicols}
}{exe:nb-antecedents}{
	\begin{enumerate}
		\item La fonction affine $f(x) = x$ donne une droite dont chaque image admet un unique antécédent.
		\item La fonction affine $f(x) = x^2$ donne une parabole. En résolvant $f(x) =1$, on peut démontrer que seuls $1$ et $-1$ sont les antécédents de $1$.
		\item Considérons $f(x) = (x-1)\cdot(x-2)\cdot(x-3)$. Résoudre $f(x)=0$ pour montrer que seuls $1 ; 2 ;$ et $3$ sont antécédents de $0$.
		\item Une fonction constante fonctionne bien. Par exemple $f(x) = 0$.
	\end{enumerate}	
}

%%%%%%%%%%%%%%%%%%%%%%%%
%%%%%%%%%% SET 6 %%%%%%%%%%
%%%%%%%%%%%%%%%%%%%%%%%% 


\exe{}{
	Compléter l'équivalence suivante.
	Soit $f$ une fonction sur un domaine $\D\subseteq\R$ et $P(x_P ;y_P)$ un point du plan tel que $x_P\in\D$.
	Alors
	\begin{align*}
		P(x_P;y_P) \in \Ccal_f && \iff && \dots\dots\dots\dots\dots\dots\dots\dots\dots
	\end{align*}
	\vspace{-20pt}
}{exe:prop-fond-enonce}{
	\begin{align*}
		(x;y) \in \Ccal_f && \iff && y = f(x)
	\end{align*}
}

\exe{}{
	Considérons la fonction $f$ sur $\R$ donnée algébriquement par
		%\[ f(x) = \frac17-x. \]
		$ f(x) = \frac17-x.$
	Pour chaque point suivant, déterminer s'il appartient à $\Ccal_f$ ou non.
	
	\begin{multicols}{3}
	\begin{enumerate}[label=\roman*)]
		\item $\left(0; \frac17\right)$
		\item $\left(-\frac17 ; 0\right)$
		\item $\left(\frac27 ; -\frac17\right)$
	\end{enumerate}
	\end{multicols}

}{exe:Cf}{
	On rappelle la propriété fondamentale
		\begin{align*}
			(x;y) \in \Ccal_f && \iff && y = f(x)
		\end{align*}
	Pour savoir si un point $(x;y)$ appartient à $\Ccal_f$, il s'agit de vérifier si l'égalité $y=f(x)$ tient.
	
	\begin{enumerate}[label=\roman*)]
		\item On applique la propriété pour $x = 0, y= \frac17$.
		D'une part, $f(x) = f(0) = \frac17 - 0 = \frac17$, et d'autre part $y=\frac17$.
		On a donc bien $y=f(x)$ pour ce couple, et il appartient à $\Ccal_f$.
			\[ \left(0; \frac17\right) \in \Ccal_f. \]
		
		\item On choisit $(x;y) = \left(-\frac17 ; 0\right)$, et on compare $f(x) = f\left(-\frac17\right) = \frac27$ à $y=0$. 
		On a donc
			\[ \left(-\frac17 ; 0\right) \not\in \Ccal_f. \]
		\item On choisit $(x;y) = \left(\frac27 ; -\frac17\right)$, et on compare $f(x) = f\left(\frac27\right) = \frac{-1}7$ à $y=-\frac17$. 
		D'où
			\[ \left(\frac27 ; -\frac17\right) \in \Ccal_f. \]
	\end{enumerate}
}


\exe{}{
	Pour chaque paire de point $P$ et de fonction $f$, déterminer si $P\in\Ccal_f$ ou si $P\not\in\Ccal_f$.
	
	%Pour chaque fonction, trouver aussi un point $Q\neq P$ qui appartient à $\Ccal_f$.
	
	\begin{multicols}{2}
	\begin{enumerate}[itemsep=5pt]
		\item $P = (0;1)$ et $f(x) = (3x- 1)^2$
		\item $P = (0;1)$ et $f(x) = (3x- 1)^3$
		\item $P = (-6;5)$ et $f(x) = \frac23 x + 7$
		\item $P = (-2;2)$ et $f(x) = 4-x$
		\item $P = (-2;2)$ et $f(x) = x^3 -2x + 8$
		\item $P = (-1;2)$ et $f(x) = x^4 - \frac1x$
		\item $P = \left(-2;\frac29\right)$ et $f(x) = \left( \frac{x}3 \right)^2 - 7$
		\item $P = \left(5 ; \sqrt{5} -3 \right)$ et $f(x) =\frac{x}{\sqrt{5}} - 3$
	\end{enumerate}
	\end{multicols}
}{exe:prop-fond}{
	\begin{enumerate}
		\item $P\in\C_f$
		\item $P\not\in\C_f$
		\item $P\not\in\C_f$
		\item $P\not\in\C_f$
		\item $P\not\in\C_f$
		\item $P\in\C_f$
		\item $P\not\in\C_f$
		\item $P\in\C_f$
	\end{enumerate}
}


\exemulticols{}{
	Donner le domaine $\D$ du repère ci-contre.
	Calculer $4$ points appartenant aux courbes $\Ccal_f$ et $\Ccal_g$ et les esquisser dans le repère.
	\setlength{\columnseprule}{1pt} 
	\begin{multicols}{2}
		$f(x) = 2x + 1$
		
		\vspace{2cm}\,
		
		\columnbreak
		
		$g(x) = 2 - x$
	\end{multicols}
	
	Quel $x\in\D$ vérifie $f(x) = g(x)$ ?
	
	\vfill\,
}{

	
	\includegraphics[page=23]{../../CM/2-figures.pdf}


}{exe:Cf-affine}{
	Un point de de la courbe représentative de $f$ prend la forme $(x ; f(x))$, où $x$ fait partie du domaine d'étude.
	
	Pour chaque fonction, les points de sa courbe représentative sont alignés.
	La courbe d'une fonction affine est en fait une droite.
	
	Les courbes représentatives sont données ci-dessous.
	\begin{center}
	\includegraphics[page=24]{../../CM/2-figures.pdf}
	\end{center}
	Graphiquement, on lit $x\approx0,3$ vérifiant $f(x) = g(x)$, mais cette solution n'est pas exacte car
		\begin{align*}
			f(0,3) = 2\cdot0,3+1 = 1,6 && \et && g(0,3) = 2-0,3 = 1,7.
		\end{align*}
	Algébriquement, on obtient la solution exacte suivante.
		\begin{align*}
			f(x) &= g(x) \\
			2x + 1 &= 2 - x \\
			3x + 1 &= 2 \\
			3x &= 1 \\
			x &= \frac13
		\end{align*}
}


\exe{, difficulty=1}{
	Soit $f(x) = (x-12)(x-13)$.
	Développer l'expression pour vérifier que $f(x) = x^2 - 25x + 156$.
	
	Calculer les images de 0 ; 1 ; 11 ; 12 ; 13; et 14 par $f$ avec la forme adéquate et sans calculatrice.
}{exe:image-selon-forme}{
	Pour $f(0)$ et $f(1)$, on préférera la forme développée $f(x) = x^2 - 25x + 156$.
	Ainsi, $f(0) = 156$ immédiatement, et $f(1) = 1 - 25 + 156 = 132$.
	
	Pour $f(11), f(12), f(13), f(14)$, la forme factorisée est préférable, car on calcule facilement que
		\begin{align*}
			f(11) &= (11-12)(11-13) = (-1)(-2) = 2 \\
			f(12) &= (12-12)(12-13) = 0 \\
			f(13) &= (13-12)(13-13) = 0 \\
			f(14) &= (14-12)(14-13) = (2)(1) =2
		\end{align*}
}


\exe{, difficulty=1}{
	Considérons les fonctions $f(x) =x^2 -6x + 9$ et $g(x)=(x-3)^2$ pour tout $x\in\R$.
	\begin{enumerate}
		\item
		Calculer les images par $f$ et $g$ de $0$ ; de $6$ ; de $3$ ; de $-\frac23$.
		\item
		Donner deux antécédents de $9$ par $f$.
		\item
		Montrer que $f(x) = g(x)$ pour tout $x\in\R$ réel.
	\end{enumerate}
}{exe:forme-canonique}{
	\begin{enumerate}
		\item
		Pour $f$, on calcule
			\begin{align*}
				f(0) &= 9, \\
				f(6) &= 6^2 - 6^2 + 9 = 9, \\
				f(3) &= 3^2 - 6\times3 + 9 = 0, \\
				f\left(-\frac23\right) &= \left(-\frac23\right)^2 - 6\left(-\frac23\right) + 9 = \frac49 + 4 + 9 = \frac49 + \frac{117}9 = \frac{121}9.
			\end{align*}
		Pour $g$, on calcule
			\begin{align*}
				g(0) &= (-3)^2 = 9, \\
				g(6) &= 3^2 = 9, \\
				g(3) &= 0^2 = 0, \\
				g\left(-\frac23\right) &= \left(-\frac23 - 3\right)^2= \left(-\frac{11}3\right)^2 = \frac{121}9.
			\end{align*}
		\item
		0 et 6 sont deux antécédents de 9 par $f$.
		\item
		On part toujours de la forme factorisée ($g$ ici) pour arriver vers la forme développée réduite ($f$ ici).
			\begin{align*}
				g(x) &= (x-3)^2 \\
					&= (x-3)(x-3) \\
					&= x(x-3) - 3(x-3) \\
					&= x^2 - 3x - 3x + 9 \\
					&= x^2 - 6x + 9
			\end{align*}
		Le développement de carré sera travaillé pendant l'étude des identités remarquables.
	\end{enumerate}
}


\exe{, difficulty=0}{
	Un fonction $f$ définie sur tout $\R$ admet le tableau de valeurs suivant.
		\begin{center}
		\def\arraystretch{1.2}
		\setlength\tabcolsep{20pt}
		\begin{tabular}{|c|c|c|c|c|}\hline
			$x$ & 0 & -2 & 1 & -1 \\ \hline
			$f(x)$ & 1 & 0 & 0 & 1 \\ \hline
		\end{tabular}
		\end{center}
	Parmis les expressions algébriques suivantes, lesquelles \underline{ne peuvent pas} correspondre à $f(x)$ ?
		\begin{multicols}{3}
		\begin{enumerate}[label=\roman*)]
			\item $1-x$
			\item $\frac{-x^2 - x + 2}2$
			\item $1+\frac{x}2$
			\item $\frac{1-x}2$
			\item $\frac{x^3 - 3x + 2}2$
			\item $\frac{-x^4 - 2x^3 + x + 2}2$
		\end{enumerate}
		\end{multicols}
}{exe:fonctions-QCM}{
		\begin{enumerate}[label=\roman*)]
			\item
			Impossible car $f(-2) = 0$.
			\item 
			Celle-ci est possible, car toutes les images correspondent au tableau.
			\item 
			Impossible car $f(1) = 0$.
			\item
			Impossible car $f(-2) = 0$.
			\item 
			Celle-ci est possible, car toutes les images correspondent au tableau.
			\item 
			Celle-ci est possible, car toutes les images correspondent au tableau.
		\end{enumerate}
}



\exe{, difficulty=0}{
	Esquisser la courbe de la fonction $f$ donnée algébriquement par
		$ f(x) = \frac3x + 1 $
	%sur le domaine $\D=[3;10]$.
	sur $[3;10]$.
}{exe:graph-droite2}{
	Il faut choisir suffisamment de $x \in [3;10]$ du domaine afin de pouvoir bien tracer l'allure de la courbe.
	Attention, celle-ci n'est pas une droite !
	
	\begin{center}
	\includegraphics[page=9]{../../CM/2-figures.pdf}
	\end{center}
}




\exemulticols{}{
	Considérons les graphes de $f$ et $g$ sur le domaine \mbox{$\D = [-2,5 ; 2,3]$} ci-contre.

	Remplir les pointillés avec le signe $<$ (inférieur strict), $>$ (supérieur strict), ou $=$.
	\begin{align*}
		f(-2) &\quad\dots\quad g(-2) \\
		g(-2) &\quad\dots\quad f(-2) \\
		f(0) &\quad\dots\quad g(0) \\
		f(1) &\quad\dots\quad g(1) \\
		g(1) &\quad\dots\quad f(1)
	\end{align*}
	Donner l'ensemble des \mbox{$x \in [-2,5 ; 2,3]$} vérifiant $f(x) = g(x)$.
}{
	\begin{center}
	\includegraphics[page=2, scale=.9]{../../CM/2-figures.pdf}
	\end{center}
}{exe:f-equals-g}{

	\begin{align*}
		f(-2) &> g(-2) \\
		g(-2) &< f(-2) \\
		f(0) &> g(0) \\
		f(1) &< g(1) \\
		g(1) &> f(1)
	\end{align*}

	Les $x$ vérifiant $f(x)=g(x)$ sont l'abscisse des points d'intersection des courbes, donc approximativement $x\approx-2,2$ ; $x\approx0,5$ et $x \approx 2,1$.
}




\exemulticols{, difficulty=1}{
	Considérons deux fonctions, $g$ et $h$, données graphiquement sur $\D = ]{-}3,4 ; 2,3[$ ci-contre.
	
	Donner approximativement l'ensemble des nombres $x$ du domaine vérifiant les (in)équations suivantes.
	\begin{enumerate}
		\itemsep1em 
		\item $g(x) = h(x)$
		\item $g(x) \leq h(x)$
		\item $h(x) \leq g(x)$
	\end{enumerate}	
}{
	\begin{center}
	\includegraphics[page=12]{../../CM/2-figures.pdf}
	\end{center}
}{exe:eq-ineq}{
	\begin{enumerate}
		\item On cherche ici l'ensemble
			\[ \bigset{ x \in \D \text{ tels que } g(x) = h(x) }. \]
		On lit (approximativement) les abscisses des points d'intersection des deux courbes :
			\[ \bigset{ x \in \D \tq g(x) = h(x) } \approx \bigset{ {-0,8}, {1,2} } \]
		\item On cherche ici l'ensemble
			\[ \bigset{ x \in \D \tq g(x) \leq h(x) }. \]
		On lit (approximativement) les abscisses pour lesquelles la courbe $\Ccal_g$ est sous $\Ccal_h$ :
			\[ \bigset{ x \in \D \tq g(x) \leq h(x) } \approx [ {-0,8} ; {1,2} ]. \]
		\item On cherche ici l'ensemble
			\[ \bigset{ x \in \D \tq h(x) \leq g(x) }. \]
		On lit (approximativement) les abscisses pour lesquelles la courbe $\Ccal_h$ est sous $\Ccal_g$.
			\[ \bigset{ x \in \D \tq g(x) \leq h(x) } \approx [{-3,4} ; {-0,8} ] \cup [ {1,2} ; {2,3} ]. \]
	\end{enumerate}
	
	Remarquons en outre les relations suivantes :
		\[ \bigset{ x \in \D \tq h(x) \leq g(x) } \cup \bigset{ x \in \D \tq h(x) \geq g(x) } = \D, \]
	ce qui n'est pas surprenant car pour tout $x\in\D$, on a forcément $g(x) \leq h(x)$ \textbf{ou} $g(x) \geq h(x)$. 
	Ainsi, 
		\[ \bigset{ x \in \D \tq h(x) \leq g(x) \ou  h(x) \geq g(x) }  = \bigset{ x \in \D } = \D. \]
	C'est donc pour cela que les ensembles sont complémentaires par rapport à $\D$ : leur union est $\D$ tout entier.
	
	En outre, les points appartenant aux deux ensembles à la fois sont
		\begin{align*}
		 \bigset{ x \in \D \tq h(x) \leq g(x) \et h(x) \geq g(x) } = \bigset{ x \in \D \tq h(x) = g(x) },
		 \end{align*}
	car $h(x) \leq g(x)$ \textbf{et} $h(x) \geq g(x)$ est équivalent à $g(x) = h(x)$ pour tout $x\in\D$.
	
	C'est donc pour ça que les bornes des intervalles étudiés sont les abscisses des points d'intersection.
	Ces abscisses sont exactement où une courbe passe au-dessus d'une autre.
}


\exe{, difficulty=1}{
	Utiliser le repère ci-dessous pour comparer les représentations graphiques des trois fonctions suivantes données algébriquement.
		\begin{align*}
			f(x) = x^2, && g(x) = x^2 - 3, && h(x) = (x+4)^2.
		\end{align*}
		
	\begin{center}
	\includegraphics[page=10]{../../CM/2-figures.pdf}
	\end{center}
	
	\begin{enumerate}[label=(\alph*)]
		\item
		Estimer, grâce aux courbes, les solutions $x\in[-8;5]$ de l'équation $f(x) = h(x)$.
		Vérifier l'exactitude des solutions en calculant et en comparant $f(x)$ et $h(x)$.
		\item
		Faire de même pour les solutions $x\in[-8;5]$ de $g(x) = h(x)$.
		\item
		Estimer, grâce aux courbes, à quel intervalle appartiennent les $x\in[-8;5]$ vérifiant \mbox{$h(x) \leq f(x)$}.
	\end{enumerate}
}{exe:y-shift}{
	On graphe les fonctions ci-dessous. 
	On remarque que $\Ccal_g$ est juste en dessous de $\Ccal_f$. C'est logique car
		\[ g(x) = f(x) - 3, \]
	donc quand on place les points pour $\Ccal_f$ et $\Ccal_g$, ceux de $\Ccal_g$ sont $3$ en dessous de ceux de $\Ccal_f$.
	
	On remarque aussi que $\Ccal_h$ est la courbe de $\Ccal_f$ décalée vers la gauche.
	C'est également cohérent car
		\[ h(x) = f(x+4). \]
	Pour connaître l'image de $x$ par $h$, on prend l'image de $x+4$ par $f$.
	On a donc par exemple $h(-4) = f(0), h(3,9) = f(0,1),$ etc... 
	On comprend bien pourquoi $\Ccal_h$ est la $\Ccal_f$ décalée de $4$ vers la gauche.
	
	\begin{center}
	\includegraphics[page=11	]{../../CM/2-figures.pdf}
	\end{center}

	\begin{enumerate}[label=(\alph*)]
		\item
		Les $x\in[-8;5]$ vérifiant $f(x)=h(x)$ sont les abscisses des points d'intersection de $\Ccal_f$ et $\Ccal_h$.
		Ici, seul $x\approx -2$ fonctionne.
		On calcule $f(-2) = 4$ et $h(-2) = 4$, et on en déduit que $x=-2$ exactement.
		
		\item
		Les $x\in[-8;5]$ vérifiant $h(x)=g(x)$ sont les abscisses des points d'intersection de $\Ccal_h$ et $\Ccal_g$.
		Ici, seul $x\approx -2,5$ fonctionne.
		On calcule $h(-2,5) = 1,5^2 = \frac94$ et $g(-2,5) = 2,5^2 - 3 = \frac{25}4 - 3 = \frac{13}4$ et donc que la solution trouvée n'est pas exacte.
		
		La solution exacte est en fait $x= -\frac{19}8 = -2,375$.
		\item
		Les $x\in[-8;5]$ vérifiant $h(x) \leq f(x)$ sont les abscisses pour lesquelles la courbe $\Ccal_h$ est sous la courbe $\Ccal_f$.
		Ici, on trouve approximativement $[-8 ; -2]$.
	\end{enumerate}
}



\exe{}{
	Calculer ou lire le coefficient directeur $a$ et l'ordonnée à l'origine $b$ des fonctions affines $f(x) = ax+b$ suivantes.
	
	\begin{multicols}{2}
	\begin{enumerate}
		\item $f(x) = 2x+1$
		\item $g(x) = 2- \frac23 x $
		\item $h(x) = x$
		\item $i(x) = 2-x$
		\item $j(x) = \frac{x}2+1$
		\item $k(x) = 7$
		\item $l(x) = 0$
		\item $m(x) = -\frac{3x}{6}$
	\end{enumerate}
	\end{multicols}
}{exe:lire-a-b}{
	
}


\exe{, difficulty=1}{
	Considérons la courbe représentative de la fonction $f(x) = x + 1$ et le point $A(2;1)$.
	Le but de l'exercice est de trouver le point de $\Ccal_f$ le plus proche de $A$.
	\begin{enumerate}
		\item
		Grapher $\Ccal_f$ sur le domaine $[0 ; 4]$ et placer le point $A$ dans un repère orthonormé.
		Le point $A$ appartient-il à $\Ccal_f$ ?
		\item
		Montrer qu'un point de $\Ccal_f$ est de la forme $(x ; x+1)$, avec $x\in\R$.
		\item
		Montrer que la distance au carré de ce point à A est égale à
			\[ 2(x-1)^2 + 2. \]
		\item
		Conclure que le point de $\Ccal_f$ le plus proche de $A$ est $B(1 ; 2)$.
		Placer $B$ dans le repère. %déjà construit.
	\end{enumerate}
}{exe:proj-xp1}{
	\begin{center}
	\includegraphics[page=3]{../../CM/2-figures.pdf}
	\end{center}
	\begin{enumerate}
		\item
		D'après la propriété fondamentale, $A \not\in \Ccal_f$ car $f(x_A) = f(2) = 2+1 = 3 \neq y_A$.
		\item
		D'après la propriété fondamentale, un point $(x ; y)$ appartient $\Ccal_f$ si et seulement si $y = f(x) = x+1$.
		Tous les points de $\Ccal_f$ sont donc de la forme $(x ; x+1)$.
		\item
		D'après la formule de la distance au carré, celle-ci est donnée par
			\begin{align*}
				(x - x_A)^2 + (x+1 - y_A)^2 &= (x-2)^2 + (x+1-1)^2,
											\\ &= (x-2)(x-2) + x^2,
											 \\ &= x^2 -2x - 2x + 4 + x^2,
											  \\ &= 2x^2 - 4x + 4.
			\end{align*}
		D'autre part, en développant $2(x-1)^2 + 2$, on obtient
			\begin{align*}
				2(x-1)^2 + 2 &= 2(x-1)(x-1) + 2,
							\\ &= 2[x^2 - x - x + 1] + 2,
							\\ &= 2x^2 - 4x + 2 + 2,
							\\ &= 2x^2 - 4x + 4,
			\end{align*}
		ce qui montre l'identité recherchée.
		\item
		Remarquons qu'un carré est toujours positif : l'expression $2(x-1)^2 + 2$ est donc toujours supérieure ou égale à 2, avec égalité quand $x-1$ est nul, c'est-à-dire quand $x=1$.
		En $x=1$, le point $(x ; x+1)$ correspond bien au $B(1 ; 2)$ proposé.
	\end{enumerate}
}


\exe{, difficulty=1}{
	Considérons la représentation graphique suivante d'une fonction $f$ définie sur \mbox{$\D = ]{-}3,4 ; 2,3[$}.
	
	\begin{center}
	\includegraphics[page=13]{../../CM/2-figures.pdf}
	\end{center}
	\begin{enumerate}
		\item\label{q1}
		Donner approximativement les images de $-1$ et de $2$ par $f$. 
		\item Énumérer approximativement les antécédents de 2 et de 4 par $f$.
		\item Exprimer aproximativement l'ensemble $\{ x \in \D \tq f(x) \geq 4 \}$ sous forme d'intervalle.
		\item Donner approximativement un réel qui admet exactement deux antécédents par $f$.
		\item Si $f$ était définie sur $\R$ tout entier, serait-il toujours possible de connaître l'image de $-2$ ? Et tous les antécédents de $-2$ ?
	\end{enumerate}
	Supposons désormais que $f(x) = 3-2 x +\frac13 x^3$ pour tout $x\in\D$ du domaine.
	\begin{enumerate}[resume]
		\item Répondre à nouveau à la question \ref{q1} à l'aide de l'expression algébrique de $f$.
		Des valeurs exactes sont attendues.
		\item Montrer que l'image de $-3$ par $f$ est $0$ et que l'image de $0$ par $f$ est $3$.
	\end{enumerate}
}{exe:deg3}{
	\begin{enumerate}
		\item On détermine approximativement
			\begin{align*}
			f(-1) \approx 4,5 && \et && f(2) \approx 1,5.
			 \end{align*}
		\item On cherche d'abord les antécédents de 2, c'est-à-dire les nombres $x$ du domaine vérifiant
			\[ f(x) = -2. \]
		Pour ça, on se pose \og à hauteur 2 \fg : on trace une droite horizontale d'ordonnée 2 et on regarde les points d'intersection.
			\begin{center}
		\includegraphics[page=14]{../../CM/2-figures.pdf}
			\end{center}
		
		
		En l'occurrence, seuls $x\approx -2,8 ; 0,5 ;  \et 2,1$ fonctionnent.
			
		Pour les antécédents de 4, on fait idem en regardant les points d'intersection avec la droite d'ordonnée 4.
			\begin{center}
		\includegraphics[page=15]{../../CM/2-figures.pdf}
			\end{center}
		On trouve trois valeurs approximatives : $x \approx -2,2$ et $x \approx 0,5$.
		
		\item
		On regarde quels antécédents $x \in \D$ du domaine d'étude ont une image supérieure ou égale à 4.
		En reprenant le graphe de la résolution de $f(x) = 4$, on remarque que \textbf{tous} les $x$ entre -2,2 et 0,5 ont une image supérieure à 4.
		Ainsi, 
			\[ \bigset{ x \in \D \tq f(x) \geq 4 } \approx [-2,2 ; 0,5], \]
		les bornes étant incluses car l'inégalité est large.
		
		\item 
		On a plusieurs choix ici. 
		Il faut placer une droite horizontale telle qu'elle s'intersecte exactement deux fois avec la courbe de $f$.
		Un choix clair est $4$ (en violet ci-dessous).
		Un choix moins clair est $1,1$, en faisant en sorte que la courbe frôle la droite horizontale qu'on place en ordonnée $1,1$ (en rouge ci-dessous).
		On dit alors que la droite est \emph{tangente} à la courbe.
		
			\begin{center}
		\includegraphics[page=16]{../../CM/2-figures.pdf}
			\end{center}
		
		\item L'image est unique et ne dépend pas du domaine (tant que celui-ci contient $-2$ !).
		On peut donc toujours connaître $f(-2)$, même sans connaître $f$ en dehors du domaine.
		
		Les antécédents, eux, dépendent du domaine choisi.
		Comme on ne sait pas du tout à quoi ressemble $\Ccal_f$ en dehors du domaine choisi, il est impossible de déterminer tous les réels $x\in\R$ antécédents de $-2$ par $f$.
		
		\item On calcule grâce à la calculatrice (ou sans...)
			\begin{align*}
				f(-1) = \frac{14}3 && f(2) =  \frac53.
			\end{align*}
		
		\item On calcule à la main que
			\begin{align*}
				f(-3) &= 3 - 2 \cdot (-3) + \frac13 \cdot (-3)^3, \\
					&= 3 + 6 - 9, \\
					&= 0.
			\end{align*}
		Pour l'image de $0$, remarquons que seul le terme ne dépendant pas de $x$ subsiste. 
		On l'appelle le terme \emph{constant}, et on trouve $f(0) = 3$, comme requis.
	\end{enumerate}
}


\exe{}{
	Tracer la courbe de la fonction $f$ sur $\D=[-5;3]$ donnée algébriquement par
		\[ f(x) = 1-x. \]
	Que dire de $\Ccal_f$ ?
	
	\emph{Rappel : $\Ccal_f$ est l'ensemble des points $\bigl( x ; f(x) \bigr)$ où $x$ parcourt le domaine $\D$.}
}{exe:graph-droite}{		
	$\Ccal_f$ est une droite.
	\begin{center}
	\includegraphics[page=1]{../../CM/2-figures.pdf}
	\end{center}
		
}




\exe{, difficulty=0}{
	Soit $f$ la fonction définie sur $\R$ par
		\[ f(x) = -10 + 3(3x-1)^2. \]
	\begin{enumerate}
		\item Développer réduire l'expression algébrique de $f$.
		\item Montrer que $f$ atteint son minimum en $x^\star=\frac13$ et donner sa valeur.
	\end{enumerate}
}{exe:extrema4-99}{
	Développons et réduisons calmement :
		\begin{align*}
			f(x) &= -10 + 3(3x-1)^2 \\
				&= -10 + 3 \bigl[ (3x)^2 + 1^2 - 2\times3x \bigr] \\
				&= -10 + 3 \bigl[ 9x^2 + 1 - 6x \bigr] \\
				&= -10 + 27x^2 + 3 - 18x \\
				&= 27x^2 - 18x - 7
		\end{align*}

	Lorsqu'on a affaire à une forme avec carré, on part systématiquement du fait qu'un carré est toujours positif pour construire $f(x)$.
	Pour tout $x\in\R$ réel, on a donc
		\begin{align*}
			(3x-1)^2 &\geq 0 \\
			3(3x-1)^2 &\geq 0 \\
			-10 + 3(3x-1)^2 &\geq -10 \\
			f(x) &\geq -10
		\end{align*}
	On en déduit que $f(x)$ est borné inférieurement par $-10$ pour tous les $x\in\R$ réels.
	
	Pour montrer que c'est un minimum atteint en $\frac13$, on calcule $f(\frac13) = -10$.
	En conclusion,
		\[ f(x) \geq f\left(\frac13\right)=-10, \]
	et ce pour tous les $x\in\R$. Par définition, $-10$ est le minimum de $f$, atteint en $\frac13$.
}

\exe{, difficulty=0}{
	Soit $f$ la fonction définie sur $\R$ par
		\[ f(x) = -3 - (x+1)^2. \]
	\begin{enumerate}
		\item Développer réduire l'expression algébrique de $f$.
		\item Donner le maximum de $f$ ainsi que l'antécédent $x^\star$ qui le réalise.
	\end{enumerate}
}{exe:extrema3-99}{
	Développons et réduisons en faisant attention au signe moins qui multiplie le carré tout entier :
		\begin{align*}
			f(x) &= -3 - (x+1)^2 \\
				&=-3 - \bigl[ x^2 + 2x + 1 \bigr] \\
				&= -3 -x^2 - 2x - 1 \\
				&= -x^2 - 2x - 4
		\end{align*}
		
	Lorsqu'on a affaire à une forme avec carré, on part systématiquement du fait qu'un carré est toujours positif pour enfin construire $f(x)$.
	Pour tout $x\in\R$ réel, on a donc
		\begin{align*}
			(x+1)^2 &\geq 0 \\
			-(x+1)^2 &\leq 0 \\
			-3 - (x+1)^2 &\leq -3 \\
			f(x) &\leq -3
		\end{align*}
	On en déduit que $f(x)$ est borné supérieurement par $-3$ pour tous les $x\in\R$ réels.
	
	Pour montrer que c'est un maximum atteint en $-1$, on calcule $f(-1) = -3 - (-1+1)^2 = -3 + 0^2 = -3$.
	En conclusion,
		\[ f(x) \leq f(-1)=-3, \]
	et ce pour tous les $x\in\R$. Par définition, $-3$ est le maximum de $f$, atteint en $-1$.
}


\exe{}{
	Montrer que la notation $x = \sqrt{-1}$ n'a pas de sens pour $x\in\R$ réel.
}{exe:sqrt-undef}{
	À supposer qu'il existe un tel $x\in\R$, alors il vérifierait $x^2 = -1$.
	Ce n'est pas possible car la fonction carré est toujours positive. \Large\Lightning
}




\exe{}{
	%Trouver le point d'intersection des droites $\Ccal_f$ et $\Ccal_g$ en posant $f(x) = g(x)$ et en résolvant pour $x$, l'abscisse du point, puis en calculant l'image $f(x)$ pour obtenir l'ordonnée du point.
	Trouver le point d'intersection des droites $\Ccal_f$ et $\Ccal_g$.

	\begin{multicols}{2}
	\begin{enumerate}[itemsep=10pt]
		\item $f(x) = 3x + 5$ et $g(x) = 2x + 7$
		\item $f(x) = 4x - 6 $ et $g(x) = 3x + 8$
		\item $f(x) = 5x + 10 $ et $g(x) = 2x + 4$
		\item $f(x) = -2x + 3 $ et $g(x) = 4x - 1$
		\item $f(x) = 7x - 5 $ et $g(x) = 5x + 9$
		\item $f(x) = 6x + 4 $ et $g(x) = 3x + 12$
		\item $f(x) = -3x + 2 $ et $g(x) = -5x - 4$
		\item $f(x) = 8x - 9 $ et $g(x) = 6x + 3$
		\item $f(x) = 2x + 11 $ et $g(x) = -3x + 5$
		\item $f(x) = -4x + 7 $ et $g(x) = -x + 10$
	\end{enumerate}
	\end{multicols}
}{exe:inter-affine2}
{
	\begin{multicols}{2}
	\begin{enumerate}[itemsep=10pt]
		\item $3x + 5 = 2x + 7 \iff x = 2$
		\item $4x - 6 = 3x + 8 \iff x = 14$
		\item $5x + 10 = 2x + 4 \iff x = -2$
		\item $-2x + 3 = 4x - 1 \iff x = \frac{2}{3}$
		\item $7x - 5 = 5x + 9 \iff x = 7$
		\item $6x + 4 = 3x + 12 \iff x = \frac{8}{3}$
		\item $-3x + 2 = -5x - 4 \iff x = -3$
		\item $8x - 9 = 6x + 3 \iff x = 6$
		\item $2x + 11 = -3x + 5 \iff x = \frac{-6}{5}$
		\item $-4x + 7 = -x + 10 \iff x = -1$
	\end{enumerate}
	\end{multicols}
}


\exe{, difficulty=1}{
	Donner graphiquement une fonction sur $\R$ non constante telle que toutes les images de $f$ admettent un nombre infini d'antécédents.
}{exe:infinite-antecedents}{
	On peut définir par exemple la fonction \emph{signe} en incluant 0 dans les positifs (traditionnellement, $\text{signe}(0) = 0$).
		\[ \text{signe}(x) = \begin{cases*} +1 & si $x \geq 0$, \\
								-1 & si $x<0$.
				\end{cases*}. \]
	Pour une fonction plus intéressante, on pourrait prendre une fonction en vague.
	La fonction sinus fonctionne bien : entrer par exemple \texttt{y=sin(x)} sur Geogebra.
}


\exemulticols{}{
	Tracer les graphes des fonctions $g(x) = f(x) + 2$ et $h(x) = f(x) - 1$ dans le repère contenant $\Ccal_f$ ci-contre.
}{
	\begin{center}
	\includegraphics[page=17]{../../CM/2-figures.pdf}
	\end{center}
}{exe:draw-plusc}{
	\begin{center}
	\includegraphics[page=18]{../../CM/2-figures.pdf}
	\end{center}
}




\exe{, difficulty=2}{
	Considérons une fonction quadratique 
		$ f(x) = ax^2 + bx + c, $
	où $a, b, c\in\R$ sont trois paramètres réels.
	Supposons qu'on connaisse deux \emph{racines} distinctes de $f$, c'est-à-dire qu'on connaisse $\alpha, \beta\in\R$ tels que $\alpha\neq\beta$ et
		\[ f(\alpha) = f(\beta) = 0. \]
	\begin{enumerate}
		\item Montrer que la fonction $g(x) = f(x) - a (x-\alpha)(x-\beta)$ est affine.
			%\[ g(x) = f(x) - a (x-\alpha)(x-\beta) \qquad \text{ pour tout } x\in\R. \]
		\item Montrer que $g$ admet deux racines distinctes.
		\item En déduire que $g$ est constamment nulle et donc que
%			\begin{align*}
%				\alpha+\beta = -\frac{b}a, && \alpha\cdot\beta = \frac{c}a, && f(x) = a (x-\alpha)(x-\beta)  \qquad \text{ pour tout } x\in\R.
%			\end{align*}
			\[ f(x) = a (x-\alpha)(x-\beta)  \qquad \text{ pour tout } x\in\R.\]
		On dit alors que $f$ se \emph{scinde} en produit affine.
	\end{enumerate}
}{exe:thm-fond-alg-2}{
	%À rendre pour possibilité de points bonus à la prochaine évaluation.
	%Ni correction ni points ne seront attribués à un travail suspect ou ne démontrant pas une bonne compréhension de l'exercice.
	
	\underline{Discussion}
	
	Il s'avère qu'un résultat plus général est atteignable : si $\alpha\in\R$ est une racine d'un polynôme\footnote{Une expression de la forme $a_n x^n + a_{n-1}x^{n-1} + \hdots + a_1 x + a_0$ dont les $a_i$ sont des coefficients fixés.} $f$ de n'importe quel degré\footnote{La plus haute puissance de $x$ : le degré de $2x^2 -3x^4 + 12$ est 4.}, alors on a nécessairement $f(x) = (x-\alpha)\cdot g(x)$ pour un polynôme $g$.
	La démonstration se base sur la division euclidienne dans les polynômes :
	pour n'importe quels polynômes $f, g$, il existe deux polynômes $q, r$ vérifiants 
		\[ f = qg + r, \]
	avec le degré de $r$ strictement inférieur au degré de $g$.
	
	Dans notre cas, il doit exister un polynôme $q$ et une constante $r$ tels que $f(x) = q(x)\cdot(x-\alpha) + r$.
	En évaluant en $x=\alpha$, on obtient que $r=0$.
	
	En répétant le processus, tout polynôme de degré $n$ admettant $n$ racines se scinde en produit affine.
	En revanche, tous les polynômes n'admettent pas de racines réelles !
	C'est le cas du polynôme $x^2 + 1$ qui, lui, ne peut pas se scinder comme $(x-\alpha)(x-\beta)$ car $\alpha$ et $\beta$ seraient alors deux racines de $x^2 + 1$.
	 Or $x^2 + 1$ ne s'annule jamais, le carré étant toujours positif !
	 
	Supposons désormais qu'on ait un ensemble de nombres contenant $\R$ ainsi qu'un nouveau « nombre », racine hypothétique de $x^2 + 1$ ; est-ce que tous les polynômes se scindent dans cet ensemble\footnote{Il faut aussi qu'on puisse ajouter et multiplier dans cet ensemble tout en restant à l'intérieur de celui-ci ; c'est la \emph{stabilité}.} ? 
	Le \emph{théorème fondamental de l'algèbre} répond positivement à cette question.
	Ce nombre hypothétique est nommé $i$, dont la construction formelle utilise, notamment, la division euclidienne des polynômes.
}



\exe{}{
	Pour chaque proposition suivante, démontrer qu'elle est \underline{toujours vraie} ou montrer qu'elle \underline{peut être fausse} avec un contre-exemple (un dessin peut suffir).
	
	$f, g, h, F$ et $G$ sont des fonctions définies sur $\R$. Un zéro est un antécédent de 0.
	\begin{enumerate}
		\item Si $f$ est affine et admet deux zéros distincts, alors $f(x) = 0$ pour tout $x\in\R$.
		\item Si $g(x) \geq 0$ pour tout $x\in\R$, alors $g$ n'admet aucun zéro.
		\item Si $h(x) > 0$ pour tout $x\in\R$, alors $h$ n'admet aucun zéro.
		\item Si $F$ n'admet aucun zéro, alors $F(x) > 0$ pour tout $x\in\R$.
		\item $G(x) = (x-3)^2 (x-5)^2$ est toujours positive est admet exactement deux zéros distincts.
	\end{enumerate}
}{exe:fonctions14}{
	\begin{enumerate}
		\item 
		C'est vrai. 
		Si $f$ est affine et admet deux racines distinctes, alors sa courbe représentative est une droite qui intersecte deux fois l'axe des abscisses en deux endroits différents.
		$\Ccal_f$ et l'axe des abscisses sont donc confondues et on a bien $f(x) = 0$ pour tout $x\in\R$.
		
		Algébriquement, on a $(x_A ; 0), (x_B ; 0) \in \Ccal_f$ pour $x_A \neq x_B$.
		La formule du coefficient directeur donne $a = 0$, et une appartenance donne $b=0$.
		
		\item 
		C'est faux : être positif ou nul n'empêche pas d'être nul.
		On pourra prend la fonction carré $f(x) = x^2$ comme contre-exemple : $f(x) \geq 0$ pour tout $x\in\R$ et pourtant $f$ admet $0$ pour racine car $f(0) = 0$.
		
		\item 
		C'est vrai : être strictement positif empêche d'être nul. 
		C'est bien la différence entre la positivité stricte ($>$) et large ($\geq$).
		
		\item 
		C'est faux : si $F$ ne s'annule jamais, elle n'est pas nécessairement toujours strictement positive car elle pourrait aussi être toujours strictement négative !
		Prenons par exemple $f(x) = -1 -x^2$ qui vérifie $f(x) \leq -1 < 0$ pour tout $x\in\R$ réel.
		
		\item 
		C'est vrai.
		Un carré est toujours positif, donc $G(x)$ est le produit de deux positifs et est positif.
		En outre, si $G(x) = 0$, alors $(x-3)^2 = 0$ ou $(x-5)^2 = 0$.
		Il suit que $x=3$ ou $x=5$, et ce sont les deux seules racines de $G$.
		
	\end{enumerate}

}



\exe{, difficulty=1}{
	Considérons les points $O(0;0)$ et $P(3 ; 1)$ et $(d)$ la médiatrice du segment $[OP]$ : ce sont tous les points $M(x;y)$ à équidistance de $O$ et $P$.
	
	Montrer que $x$ et $y$ sont liés par la relation 
		$ y = -3x + 5. $
}{exe:mediatrice}{
%	À rendre pour possibilité de points bonus à la prochaine évaluation.
%	Ni correction ni points ne seront attribués à un travail suspect ou ne démontrant pas une bonne compréhension de l'exercice.
	La contrainte d'équidistance s'exprime et est équivalente aux équations suivantes.
		\begin{align*}
			\norm{M-O}^2 &= \norm{M - P}^2 \\
			x^2 + y^2 &= (x-3)^2 + (y-1)^2 \\
			x^2 + y^2 &= x^2 - 6x + 9 + y^2 - 2y + 1 \\
			2y &= -6x + 10 \\
			y &= -3x + 5
		\end{align*}
}


\exe{}{
	%À l'aide de l'exercice \ref{exe:1}, donner le domaine de définition de chaque fonction suivante.
	Calculer le domaine de définition de chaque fonction suivante.
	
	\begin{multicols}{3}
	\begin{enumerate}
		\item $f(x) = \sqrt{x+3}$
		\item $g(x) = \sqrt{3x-10}$
		\item $h(x) = \sqrt{-4x-12}$
	\end{enumerate}
	\end{multicols}

}{exe:domaine-def}{
	Le domaine de définition est le plus grand domaine sur lequel une fonction est bien définie.
	Ici, la seule chose qui peut se mal passer est qu'on demande à $f$ de calculer une racine carrée d'un nombre négatif.
	Ceci n'est pas possible car $x^2 \geq 0$ pour tout $x\in\R$, donc si on souhaite calculer $\sqrt{a}$ vérifiant $\sqrt{a}^2 = a$ par définition, on a nécessairement $a\geq0$.
	Par exemple, $\sqrt{-1}$ ne peut pas être un nombre réel, sinon $\sqrt{-1}^2 = -1$, et un carré serait négatif !
	
	On impose donc que l'expression sous une racine carrée soit positive ou nulle, et on prend tous les réels vérifiant cette propriété.
	\begin{enumerate}
		\item
			\[ \D_f = \{ x \in \R \text{ tq. } x + 3 \geq 0 \} = [-3; \pinfty[. \]
		\item
			\[ \D_f = \{ x \in \R \text{ tq. } 3x- 10 \geq 0 \} = \left[ \frac{10}3 ; \pinfty \right[. \]
		\item
			\[ \D_f = \{ x \in \R \text{ tq. } -4x-12 \geq 0 \} = ]\minfty; -3]. \]
	\end{enumerate}
}


\exemulticols{}{
	Considérons la fonction $h$ donnée graphiquement ci-contre.
	Sans davantage d'informations, son domaine de définition est \mbox{$\D_h = [-3;3]$}.
	
	\begin{enumerate}
		\item Donner le domaine de définition de $\frac{2,5}{h(x)}$.
		\item Donner le domaine de définition de $\frac{h(x)^2}{3}$.
		\item Donner le domaine de définition de $\frac{h(x)}{x}$.
	\end{enumerate}
}{
	\begin{center}
	\includegraphics[page=25]{../../CM/2-figures.pdf}
	\end{center}

}{exe:domaine-def-graphique}{
	\begin{enumerate}
		\item 
		Il s'agit de vérifier que le dénominateur ne s'annule pas.
		Les valeurs $-2$ et 2 sont donc interdites à $\frac{2,5}{h(x)}$ et son domaine de définition est
			\[ [-3 ; -2 [ \cup ]-2 ; 2[ \cup ]2 ; 3]. \]
		\item 
		Rien n'est interdit ici : on ne divise jamais par zéro, et on ne prend aucune racine carrée de nombre négatif.
		Le domaine de définition est donc $[-3 ; 3]$.
		\item 
		Il s'agit de vérifier que le dénominateur ne s'annule pas.
		La valeur 0 est donc interdite à $\frac{h(x)}{x}$ et son domaine de définition est
			\[ [-3 ; 0 [ \cup ]0 ; 3]. \]
		
	\end{enumerate}
}

\exe{}{
	Quel est le domaine de définition de la fonction $\xi$ suivante ?
		\[\xi(x) = \frac{42}{1+x^2}\]
}{exe:sqr-p-1}{
	Il s'agit de vérifier que le dénominateur ne s'annule pas.
	Or l'égalité $1+x^2 = 0$ est équivalente à $x^2 = -1$ qui n'admet aucune solution réelle car le carré est toujours positif.
	Par conséquent, \mbox{$\D_\xi = \R$.}
}


\exe{}{
	Calculer le domaine de définition de chaque fonction suivante.
	
	\begin{multicols}{3}
	\begin{enumerate}
		\item $F(x) = \ln(x^2 + 2x+1)$
		\item $G(x) = \frac1{2x+3}-\ln(2x+6)$
		\item $H(x) = \sin\left(\frac1{2x}\right)$
	\end{enumerate}
	\end{multicols}

}{exe:domaine-def-log}{
	\begin{enumerate}
		\item
		Il faut imposer $x^2 + 2x+1 > 0$.
		Or, comme $x^2 + 2x + 1 = (x+1)^2$, cette quantité est strictement positive pour tout $x\in\R$ sauf $1$.
		Par conséquent, $\D_F = \R-\{1\}$.
		\item 
		Il faut imposer que $2x+3\neq0$ et que $2x+6>0$.
		L'équation donne $x\neq-\frac32$ et l'inéquation donne $x>-3$.
		D'où $\D_G = ]-3 ; \pinfty[ - \{-\frac32\} = ]-3 ; -\frac32[ \cup ]-\frac32 ; \pinfty[$.
		\item
		La fonction sinus n'a aucune contrainte particulière, mais la fonction inverse nous impose que $2x \neq 0 \iff x \neq 0$.
		Donc $\D_H = \R^* = \R - \{ 0 \}$.
	\end{enumerate}
}



\exe{,difficulty=2}{
	Soit $f(x) = x^2 - x - 1$ une fonction quadratique sur $\R$.
	
	\begin{enumerate}
		\item
		Montrer que  $f(x) = \left( x - \frac{1 - \sqrt{5}}2 \right)\cdot \left( x - \frac{1 + \sqrt{5}}2 \right)$ pour tout $x\in\R$.
		\item 
		En déduire les zéros de $f$. Lequel est positif et lequel est négatif ? Raisonner sans calculatrice.
	\end{enumerate}
	%Le zéro positif $\phi = \frac{1 + \sqrt{5}}2$ (lu \og phi \fg) s'appelle le \emph{nombre d'or}.
	Le zéro positif $\phi$ (lu \og phi \fg) s'appelle le \emph{nombre d'or}.
	\begin{enumerate}[resume]
		\item
		Vérifier que les zéros trouvés annulent bien $f$ en calculant leur image par $f$.
	\end{enumerate}
}{exe:phi}{
%	À rendre pour possibilité de points bonus à la prochaine évaluation.
%	Ni correction ni points ne seront attribués à un travail suspect ou ne démontrant pas une bonne compréhension de l'exercice.
}

\exe{, difficulty=2}{
	Considérons l'équation quadratique $3x^2-7x+4= 0$ pour $x\in\R$.
	\begin{enumerate}
		\item
		Montrer que $1$ est racine de l'équation.
	\end{enumerate}
	Le but des questions suivantes est d'effectuer la division euclidienne de $3x^2-7x+4$ par $x-1$ afin de trouver la deuxième racine de l'équation.
	\begin{enumerate}[resume]
		\item
		Quelle fonction affine $q_1(x)$ prendre pour que $3x^2 -7x+4 - (x-1)q_1(x)$ n'ait plus de terme en $x^2$ ?
		\item
		Quelle fonction affine $q_2(x)$ prendre pour que $3x^2 -7x+4 - (x-1)q_1(x) - (x-1)q_2(x)$ n'ait plus de terme en $x$ ?
		\item
		Montrer que $3x^2 - 7x+4 = (x-1)\bigl[q_1(x) + q_2(x)\bigr]$ et en déduire l'autre racine de l'équation.
	\end{enumerate}
}{exe:euclide-poly}{
%	À rendre pour possibilité de points bonus à la prochaine évaluation.
%	Ni correction ni points ne seront attribués à un travail suspect ou ne démontrant pas une bonne compréhension de l'exercice.
}

%c'est dur en vrai
\exe{, difficulty=2}{
	Soient $m, n, d\in\Z$ trois entiers relatifs.
	Donner explicitement une fonction $f(x) = ax^2 + bx+c$ avec $a, b, c$ entiers relatifs
	telle que $f\bigl(m+n\sqrt{d}\bigr) = 0$.
	
	En déduire une fonction à coefficient entiers qui annule $1-\sqrt2$.
}{exe:all-roots}{
%	etc...
}

\exe{,difficulty=2}{
	Soit $f(x) = ax^2 + bx + c$ un polynôme à coefficients $a, b, c$ entiers.
	Soit $d\in\N$ un entier naturel qui n'est pas un carré parfait.
	\begin{enumerate}
		\item Montrer que si $f\bigl(\sqrt{d}\bigr)= 0$, alors $b=0$. Raisonner par contradiction.
		\item En déduire que si $\sqrt{d}$ est racine de $f$, alors $-\sqrt{d}$ est nécessairement l'autre racine de $f$.
		\item Une racine du polynôme $f(x) = x^2 - x - 1$ peut-elle s'écrire $\sqrt{k}$ avec $k\in\N$ un entier non carré parfait ?
		%\item Donner des paramètres $a,b,c$ entiers avec $b\neq0$ tels que $f$ admette une racine entière.
	\end{enumerate}
}{exe:quadratic-surd}{
%	etc...
}

\exe{, difficulty=2}{
	Supposons que $r+s\sqrt{d}$, avec $r, s\in\Q$ et $d\in\N$ non carré parfait, soit racine de l'équation $ax^2 + bx + c = 0$ où $a, b, c\in\Q$.
	Montrer que son conjugué $r-s\sqrt{d}$ est également racine.
}{exe:Galois2}{
%	À rendre pour possibilité de points bonus à la prochaine évaluation.
%	Ni correction ni points ne seront attribués à un travail suspect ou ne démontrant pas une bonne compréhension de l'exercice.
	
	Voici une première stratégie de preuve possible.
	\begin{enumerate}
		\item
		Montrer que $1$ et $\sqrt{d}$ sont \emph{linéairement indépendants} sur $\Q$.
		C'est-à-dire, montrer que si $r + s\sqrt{d} = 0$ avec $r, s \in\Q$, alors $r=s=0$.
		Une équation sur les réels devient alors deux équations sur les rationnels !
		\item
		En déduire que $a(r+s\sqrt{d})^2 + b(r+s\sqrt{d}) + c=0$ est équivalente aux deux équations $ar^2 + ads^2+br+c=0$ et  $2ars+bs=0$.
		\item
		Conclure en calculant $a(r-s\sqrt{d})^2 + b(r-s\sqrt{d}) + c$.
	\end{enumerate}
	À l'étape 2, nous utilisons aussi que la somme et le produit de deux rationnels est un rationnel, et c'est à ce moment que nous utilisons que tous les coefficients ($r, s, a, b, c$) sont rationnels.
	Cette propriété de $\Q$ se démontre assez simplement avec les formules $\frac{a}b+\frac{c}d$ et $\frac{a}b\times\frac{c}d$.
	
	Voici un deuxième stratégie de preuve possible.
	\begin{enumerate}
		\item
		Montrer que $1$ et $\sqrt{d}$ sont \emph{linéairement indépendants} sur $\Q$.
		C'est-à-dire, montrer que si $r + s\sqrt{d} = 0$ avec $r, s \in\Q$, alors $r=s=0$.
		\item
		Montrer que la fonction « conjugué » $f\bigl(r+s\sqrt{d}\bigr) = r - s\sqrt{d}$ où $r, s \in \Q$ vérifie que
			\begin{enumerate}[label=(\roman*)]
				\item $f(x+y) = f(x)+f(y)$
				\item $f(xy) = f(x)f(y)$
			\end{enumerate}
		\item
		Conclure en appliquant $f$ à l'équation $ax^2 +bx + c = 0$ pour obtenir $af(x)^2 + bf(x) + c = 0$.
	\end{enumerate}
	
}




\exe{, difficulty=3}{
	Posons $\Qsqrt{5} = \bigset{ r + s \sqrt5 \tq r, s \in \Q} \subseteq \R$ et considérons $x = r + s\sqrt5$ et $y=p+q\sqrt5$ avec $r, s, p, q \in\Q$.
	\begin{enumerate}
		\item
		Montrer que $\Q \subseteq \Qsqrt5$ est une inclusion stricte.
		\item
		Montrer que $x=y \iff r = p \et s = q$.
	\end{enumerate}
	Considérons désormais la fonction « conjugué » $f$ définie par $f\bigl(r + s\sqrt5\bigr) = r - s\sqrt5$.
	\begin{enumerate}[resume]
		\item
		Montrer que $f(x+y) = f(x)+f(y)$ et que $f(xy) = f(x)f(y)$.
		\item
		Montrer que si $f(x) = x$, alors $x\in\Q$.
		\item
		Conclure à l'aide des deux dernières questions que
			\[ \frac1{2^{256}\sqrt5}\left[ \left({1+\sqrt5}\right)^{256} - \left({1-\sqrt5}\right)^{256}\right] \]
		est un nombre rationnel.
		(C'est en fait un nombre entier et le terme $256$ de la suite de Fibonacci)
		\item
		Montrer que $\Qsqrt{5} \neq \R$ en montrant que $\sqrt2 \not\in \Qsqrt{5}$.
	\end{enumerate}
} {exe:Qsqrt2}{
	%À rendre pour possibilité de points bonus à la prochaine évaluation.
	%Ni correction ni points ne seront attribués à un travail suspect ou ne démontrant pas une bonne compréhension de l'exercice.
	
	Indication pour la question 2 : il est possible de montrer en premier lieu que $x = 0 \iff r=s=0$.
	On dit alors que $1$ et $\sqrt5$ sont \emph{linéairement indépendants} sur $\Q$.
	
	Indication pour la question 5 : montrer que $f(x^n) = f(x)^n$ et que $f\left(\frac1x\right) = \frac1{f(x)}$ à l'aide de la question 3.
	
	\underline{Discussion}
	
	L'exercice fait l'étude du corps d'extension des nombres rationnels auxquels on a ajouté la racine carrée de 5.
	Il suit de la question 6 que l'ensemble 
		\[ \Q\bigl(\sqrt2, \sqrt5\bigr) = \bigset{ r + s \sqrt2 + t\sqrt5 \tq r, s, t \in \Q} \]
	est un sur-ensemble strict de $\Qsqrt5$ qui n'est toujours pas égal à $\R$.
	
	La {théorie de Galois}\footnotemark généralise ces exemples et permet notamment de répondre par la négative aux trois grands problèmes de l'Antiquité : la duplication du cube, la trisection de l'angle, et la quadrature du cercle.
	\footnotetext{Évariste Galois (1811-1832), mathématicien français mort à l'âge de 20 ans. }
	La démonstration procède ainsi: chaque construction à la règle et au compas à partir d'un repère orthonormé ajoute la racine carrée d'un nombre déjà constructible jusqu'ici.
	La racine carrée de 2 est par exemple constructible grâce au théorème de Pythagore comme distance entre $(1 ; 0)$ et $(0;1)$.
	En montrant que certains nombres ne peuvent pas appartenir à de telles extensions, on démontre immédiatement qu'ils sont non constructibles.
}


\exe{, difficulty=1}{
	Pour chaque paire $(f, g)$ de fonctions affines sur $\R$, calculer l'intersection $\Ccal_f \cap \Ccal_g$.
	En ensemble de point(s) est attendu.
	
	\begin{multicols}{2}
	\begin{enumerate}[itemsep=10pt]
		\item $f(x) = 2x + 1$ et $g(x) = -x+1$.
		\item $f(x) = -x + 1$ et $g(x) = 2x + 1$.
		\item $f(x) = 3+7x$ et $g(x) = 2$.
		\item $f(x) = 9-2x$ et $g(x) = 17-x$.
		\item $f(x) = 2x+1$ et $g(x) = 2x+1$.
		\item $f(x) = 2x+1$ et $g(x) = 2x+2$.
	\end{enumerate}
	\end{multicols}
}{exe:affine-intersection}{

	\begin{enumerate}
		\item $f(x) = 2x + 1, g(x) = -x+1$ .
			On cherche un point $P(x_P, y_P)$ vérifiant les deux équations
				\begin{align*}
					y_P = f(x_P) && \text{ et } && y_P = g(x_P)
				\end{align*}
			Comme bien sûr $y_P = y_P$, on peut résoudre l'équation
				\[ f(x_P) = g(x_P) \]
			pour trouver $x_P$.
			Ici, on pose calmement
				\begin{align*}
					f(x_P) &= g(x_P) \\
					2x_P + 1 &= -x_P + 1 \\
					3x_P &= 0 \\
					x_P &= 0,
				\end{align*}
			ce qui nous fournit $x_P = 0$.
			Reste plus qu'à utiliser que $y_P = f(x_P) = g(x_P)$ pour calculer $y_P$.
			
			D'une part, 
				\[ f(x_P) = f(0) = 2\cdot0 + 1 = 1. \]
			D'autre part, on aurait pû calculer
				\[ g(x_P) = g(0) = -0 + 1 = 1. \]
			Rien de surprenant à ce qu'on trouve la même valeur $y_P = 1$, car $x_P$ a été trouvé vérifiant $f(x_P) = g(x_P)$.
			
			En conclusion, $\Ccal_f \cap \Ccal_g = \{ (0;1) \}$.
			
		\item 
			Similairement,
			\begin{align*}
				f(x_P) &= g(x_P) \\
				-x_P + 1 & = 2x_P + 1 \\
				x_P &= 0,
			\end{align*}
			et $y_P = f(0) = 1$.
			Sans surprise, $(0;1)$ est à nouveau le point d'intersection de $\Ccal_f$ et de $\Ccal_g$ car se sont les mêmes fonctions que ci-dessus mais échangées.
		
		\item 
			Similairement,
			\begin{align*}
				f(x_P) &= g(x_P) \\
				3 + 7x_P & = 2 \\
				7x_P &= -1 \\
				x_P &= \dfrac{-1}7,
			\end{align*}
			\[
				y_P = f(x_P) = 3+7\cdot \dfrac{-1}7  = 2.
			\]
			Il aurait été encore plus facile de calculer $g(x_P) = 2$ à la place !
			
			En conclusion, $\Ccal_f \cap \Ccal_g = \left\{ \left(-\dfrac72 ; 2\right) \right\}$.
		
		\item
			Similairement,
			\begin{align*}
				f(x_P) &= g(x_P) \\
				9-2x_P & = 17-x_P \\
				-8 &= x_P,
			\end{align*}
			\[
				y_P = f(x_P) = 9-2 \cdot(-8) = 25.
			\]
		
			En conclusion, $\Ccal_f \cap \Ccal_g = \left\{ (-8;25) \right\}$.
		
		\item 			
		Similairement,
			\begin{align*}
				f(x_P) &= g(x_P) \\
				2x_P + 1 & = 2x_P + 1 \\
				0 &= 0,
			\end{align*}
			cette équation étant évidemment vraie, tous les $x_P \in \R$ vérifient l'équation.
			\[
				y_P = f(x_P) = 2x_P + 1.
			\]
			N'importe quel $(x_P ; 2x_P+1) \in \R$ appartient donc à l'intersection des droites.
			
			En conclusion, $\Ccal_f \cap \Ccal_g = \bigset{ (x_P ; 2x_P+1) \text{ tq. } x_P \in \R } = \bigset{ (x_P ; y_P) \text{ tq. } y_P = 2x_P + 1 } = \Ccal_f = \Ccal_g$.
			Il n'est pas surprenant qu'on trouve que toute la droite appartienne à l'intersection, car $f$ et $g$ sont identiques : tous les points de la droite $\Ccal_f$ sont des points d'intersection.
		
		
		\item 
			Similairement,
			\begin{align*}
				f(x_P) &= g(x_P) \\
				2x_P + 1 & = 2x_P + 2 \\
				1 &= 2,
			\end{align*}
			cette équation étant évidemment fausse, aucun $x_P \in \R$ ne vérifie l'équation.
			
			En conclusion, $\Ccal_f \cap \Ccal_g = \emptyset$, l'ensemble vide.
			En fait, les droites $\Ccal_f$ et $\Ccal_g$ sont parallèles car elles ont le même coefficient directeur $2$.
			Or les droites ne sont pas égales car leurs ordonnées à l'origine sont différentes, donc il n'y a aucun point d'intersection.
	\end{enumerate}

}




%\ExerciseStopSelect

%%%%%%%%%%%

\newpage
\fancyhead[C]{\textbf{Solutions}}
\shipoutAnswer

\end{document}
